\documentclass{article}
\usepackage[margin=0.8in]{geometry}
\usepackage{amsmath,amssymb,amsthm,mathtools}
\usepackage{mathrsfs,bm,bbm}
\usepackage{graphicx,booktabs,array,multirow,tabularx,longtable}
\usepackage{enumitem}
\usepackage{xcolor}
\usepackage{algorithm,algpseudocode}
\usepackage{subcaption,tikz}
\usepackage{siunitx,cancel,accents}
\usepackage{microtype}
\usepackage[numbers,sort&compress]{natbib}
\usepackage[colorlinks=true,linkcolor=blue,citecolor=blue,urlcolor=blue]{hyperref}
\usepackage{cleveref}
\setlength{\emergencystretch}{2em}
\newtheorem{assumption}{Assumption}
\newtheorem{definition}{Definition}
\newtheorem{theorem}{Theorem}
\newtheorem{lemma}{Lemma}
\newtheorem{proposition}{Proposition}
\newtheorem{openendedquestion}{Open-ended Question}
\newtheorem{conjecture}{Conjecture}
\newcommand{\coreref}[1]{\ref{#1}}

\begin{document}

\sloppy

\section*{stat\_\allowbreak{}gene\_\allowbreak{}pairwise\_\allowbreak{}feasible\_\allowbreak{}minimax}

\noindent\textit{Specialization:} v1\par

\paragraph{TL;DR.} For \(d\) supplied independent cause--effect pairs, the orientation-permuted Hermite chart fills a local rectangle of GENE-inspired gains and fixed-Gaussian-kernel residual-HSIC margins. On the explicit sub-half chart \(\mathcal M_d^{\mathrm H}(\lambda)\), the observable antisymmetric moment and a median-of-means rule yield the unconditional matched all-procedure frontier \(N_{\mathrm H}^*(d,\alpha,\lambda)\asymp\lambda^{-1}\log((d+1)/\alpha)\) and a computable \(O(dn)\) simultaneous confidence set with the same singleton transition. The product-testing lower bound, but not the chart estimator or confidence set, transfers to the larger separated restricted-ANM class. Separately, one fixed chart law makes the compulsory regression-good event in the prescribed four-fold local-polynomial/orthogonalization probability guarantee fail with limiting probability above \(2/e\); this is only a fixed-protocol obstruction.

\section{Project justification}

\paragraph{Gap.} Published restricted-ANM theory gives qualitative identification, and regression--independence procedures have pointwise consistency guarantees, but no cited result gives a finite-sample matched exact-recovery frontier or a confidence-set singleton transition for many supplied independent nonlinear cause--effect pairs. Existing observed-variable HSIC testing rates do not determine the geometry or product orientation experiment considered here.

\paragraph{Niche.} Independent supplied pairs isolate two tractable objects: the feasible geometry of the GENE-inspired gain and residual-HSIC map, and an explicit sub-half Hermite-chart experiment on which one observable moment carries the causal orientation with signal squared proportional to \(\lambda\).

\paragraph{Fill.} The note proves the orientation-indexed attainable rectangle, constructs a \(d\)-query median-of-means chart rule and product confidence set, and matches their \(\lambda^{-1}\log((d+1)/\alpha)\) sufficient scale to a product-testing converse over the same chart class. Class containment transfers only the converse to the larger separated restricted-ANM class. A separate theorem gives an explicit fixed-law tail-hole witness for which the prescribed four-fold local-polynomial regression-good event has complement probability asymptotically at least \(e^{-\delta}>2/e\); because that event is a compulsory member of the fixed orthogonalization protocol's failure union, the exact protocol probability guarantee cannot hold. No broader no-statistic or full-class upper-bound claim follows.

\section{Related work}

Li et al. (2025, Equation (6) and Theorem 4.3) motivate the hybrid architecture, but their published GENE score uses neural-network explained variance and a discretized chi-square independence indicator; the oracle fixed-HSIC population score here is only GENE-inspired. Peters et al. (2014, Definition 27 and Theorem 28) supply the restricted positive-density identification domain, and their disjoint-edge specialization is used as a qualified identification and class-containment bridge, not to transfer the chart estimator or upper bound. Mooij et al. (2016, Appendix A) prove pointwise consistency of regression plus fixed-kernel independence testing under a suitable learner, whereas the present matched finite-sample frontier is proved on the explicit Hermite chart by an observable moment. Hl\'avka, Hu\v{s}kov\'a, and Meintanis (2011) explain why estimated residuals need nuisance-specific analysis; consistently, the separate theorem here refutes only the exact high-probability requirement that includes the prescribed local-polynomial event \(\mathcal G_n\). Shekhar, Kim, and Ramdas (2023) and Kim, Balakrishnan, and Wasserman (2022) provide observed-variable minimax testing templates, while Pfister et al. (2018) provide asymptotic residual-ordering arguments without this local uniform product frontier. Wang, Kolar, and Drton (2026) motivate inversion into causal-order confidence sets, and Gao and Stoev (2020) motivate the logarithmic dimension cost of exact simultaneous recovery. Acharya et al.'s discrete intervention model and Jin's linear-LiNGAM primitive-margin model are distinct experiments. Wornbard et al. receive non-archival credit for the presolve insight; no public theorem under that attribution is used as a proof premise.

\section{Setup}

Target (estimand / structural parameter): \(\theta(P)=(\theta_1(P),\ldots,\theta_d(P))\in\{-1,+1\}^d\).
Primary functional / estimator / diagnostic: \(\theta_j(P)=\operatorname*{arg\,min}_{b\in\{-1,+1\}}H_{j,b}(P)\) on the identified restricted-ANM domain. On the explicit sub-half Hermite chart, the proved observable relation \(\operatorname{sign}(\mathbb E_P\psi(X_j,Y_j))=\theta_j(P)\) yields \(N_{\mathrm H}^*(d,\alpha,\lambda)\asymp\lambda^{-1}\log((d+1)/\alpha)\) unconditionally. For the larger class, only \(N^*(d,\alpha,\lambda)\gtrsim\lambda^{-1}\log((d+1)/\alpha)\) is asserted..
\begin{itemize}
  \item \texttt{d} --- \texttt{integer}; \(\{1,2,\ldots\}\); \texttt{dimension}
  \par\smallskip\noindent\emph{Definition.} \texttt{Number of supplied cause-\allowbreak{}-\allowbreak{}effect pairs.}
  \item \texttt{I} --- \texttt{finite index set}; \texttt{pair index set}
  \par\smallskip\noindent\emph{Definition.} \(I=\{1,\ldots,d\}\).
  \item \texttt{n} --- \texttt{integer}; \(\{1,2,\ldots\}\); \texttt{sample size}
  \par\smallskip\noindent\emph{Definition.} \texttt{Number of independent observations of all supplied pairs.}
  \item \texttt{alpha} --- \texttt{scalar}; \((0,1/4]\); \texttt{confidence parameter}
  \par\smallskip\noindent\emph{Definition.} \texttt{Familywise error level.}
  \item \texttt{lambda} --- \texttt{scalar}; \((0,1]\); \texttt{signal parameter}
  \par\smallskip\noindent\emph{Definition.} \texttt{Smallest wrong-\allowbreak{}direction population residual-\allowbreak{}HSIC margin.}
  \item \texttt{beta} --- \texttt{scalar}; \((1/2,\infty)\); \texttt{regularity index}
  \par\smallskip\noindent\emph{Definition.} \texttt{Fixed Holder smoothness index of the one-\allowbreak{}dimensional mechanisms and log densities.}
  \item \texttt{L} --- \texttt{scalar}; \((1,\infty)\); \texttt{regularity constant}
  \par\smallskip\noindent\emph{Definition.} \texttt{Fixed radius for the smoothness envelope.}
  \item \texttt{c0} --- \texttt{scalar}; \((0,1)\); \texttt{regularity constant}
  \par\smallskip\noindent\emph{Definition.} \texttt{Fixed lower Gaussian-\allowbreak{}envelope constant.}
  \item \texttt{C0} --- \texttt{scalar}; \((1,\infty)\); \texttt{regularity constant}
  \par\smallskip\noindent\emph{Definition.} \texttt{Fixed upper Gaussian-\allowbreak{}envelope constant.}
  \item \texttt{sigmaK} --- \texttt{scalar}; \((0,\infty)\); \texttt{kernel constant}
  \par\smallskip\noindent\emph{Definition.} \texttt{Fixed bandwidth used by every Gaussian kernel.}
  \item \texttt{glo} --- \texttt{scalar}; \((0,1)\); \texttt{signal constant}
  \par\smallskip\noindent\emph{Definition.} \texttt{Fixed lower GENE correcting-\allowbreak{}gain margin.}
  \item \texttt{ghi} --- \texttt{scalar}; \((0,1)\); \texttt{signal constant}
  \par\smallskip\noindent\emph{Definition.} Fixed upper GENE correcting-gain margin with \(\underline g<\overline g\).
  \item \texttt{AH} --- \texttt{scalar}; \((1,\infty)\); \texttt{signal constant}
  \par\smallskip\noindent\emph{Definition.} \texttt{Fixed width factor for the local residual-\allowbreak{}HSIC window.}
  \item \texttt{lambda0} --- \texttt{scalar}; \((0,1]\); \texttt{geometry constant}
  \par\smallskip\noindent\emph{Definition.} \(\lambda_0\) is the local signal radius certified by Theorem \(\mathrm{thm:attainable\text{-}margin\text{-}rectangle}\).
  \item \texttt{theta} --- \texttt{orientation vector}; \(\{-1,+1\}^{I}\); \texttt{causal target}
  \par\smallskip\noindent\emph{Definition.} \(\theta_j=+1\) means \(X_j\to Y_j\), and \(\theta_j=-1\) means \(Y_j\to X_j\).
  \item \texttt{W} --- \texttt{random vector}; \(\mathbb R^{2d}\); \texttt{observed data unit}
  \par\smallskip\noindent\emph{Definition.} \(W=((X_j,Y_j):j\in I)\).
  \item \texttt{X} --- \texttt{random vector}; \(\mathbb R^d\); \texttt{observed coordinate}
  \par\smallskip\noindent\emph{Definition.} \(X=(X_j:j\in I)\), the first labelled coordinate of each supplied pair.
  \item \texttt{Y} --- \texttt{random vector}; \(\mathbb R^d\); \texttt{observed coordinate}
  \par\smallskip\noindent\emph{Definition.} \(Y=(Y_j:j\in I)\), the second labelled coordinate of each supplied pair.
  \item \texttt{Sample} --- \texttt{sample}; \((\mathbb R^{2d})^n\); \texttt{observed sample}
  \par\smallskip\noindent\emph{Definition.} \(\mathcal W_n=(W_1,\ldots,W_n)\) is an independent sample from \(P\).
  \item \texttt{P} --- \texttt{probability law}; \(\mathcal P(\mathbb R^{2d})\); \texttt{observed-\allowbreak{}data law}
  \par\smallskip\noindent\emph{Definition.} \(P=\mathcal L(W)\).
  \item \texttt{Ppair} --- \texttt{pair-\allowbreak{}law family}; \texttt{product factors}
  \par\smallskip\noindent\emph{Definition.} \((P_j:j\in I)\), where \(P_j=\mathcal L(X_j,Y_j)\).
  \item \texttt{Cb} --- \texttt{candidate-\allowbreak{}cause family}; \texttt{direction-\allowbreak{}indexed predictor}
  \par\smallskip\noindent\emph{Definition.} \(C_{j,+}=X_j\) and \(C_{j,-}=Y_j\).
  \item \texttt{Eb} --- \texttt{candidate-\allowbreak{}effect family}; \texttt{direction-\allowbreak{}indexed response}
  \par\smallskip\noindent\emph{Definition.} \(E_{j,+}=Y_j\) and \(E_{j,-}=X_j\).
  \item \texttt{Ctrue} --- \texttt{true-\allowbreak{}cause vector}; \texttt{causal predictor}
  \par\smallskip\noindent\emph{Definition.} \(C_j=C_{j,\theta_j}\).
  \item \texttt{Etrue} --- \texttt{true-\allowbreak{}effect vector}; \texttt{causal response}
  \par\smallskip\noindent\emph{Definition.} \(E_j=E_{j,\theta_j}\).
  \item \texttt{f} --- \texttt{mechanism family}; \(f_j:\mathbb R\to\mathbb R\); \texttt{causal primitive}
  \par\smallskip\noindent\emph{Definition.} \texttt{Structural regression functions for the true directions.}
  \item \texttt{eps} --- \texttt{noise vector}; \(\mathbb R^d\); \texttt{causal primitive}
  \par\smallskip\noindent\emph{Definition.} \(\varepsilon=(\varepsilon_j:j\in I)\).
  \item \texttt{pC} --- \texttt{density family}; \(p^C_j:\mathbb R\to(0,\infty)\); \texttt{primitive density}
  \par\smallskip\noindent\emph{Definition.} \texttt{Densities of the true causes.}
  \item \texttt{pEps} --- \texttt{density family}; \(p^\varepsilon_j:\mathbb R\to(0,\infty)\); \texttt{primitive density}
  \par\smallskip\noindent\emph{Definition.} \texttt{Densities of the structural noises.}
  \item \texttt{Hw} --- \texttt{function class}; \texttt{fixed smoothness envelope}
  \par\smallskip\noindent\emph{Definition.} \(\mathcal H_{(2)}^\beta(L)=\{u:\max_{0\le r\le\lfloor\beta\rfloor}\sup_{z\in\mathbb R}(1+|z|)^{-2}|u^{(r)}(z)|+\sup_{x\ne y}\frac{|u^{(\lfloor\beta\rfloor)}(x)-u^{(\lfloor\beta\rfloor)}(y)|}{|x-y|^{\beta-\lfloor\beta\rfloor}(1+|x|+|y|)^2}\le L\}\).
  \item \texttt{Mbase} --- \texttt{model class}; \texttt{fixed positive-\allowbreak{}density restricted-\allowbreak{}ANM envelope}
  \par\smallskip\noindent\emph{Definition.} \(\mathcal M_d(\beta,L,c_0,C_0)\).
  \item \texttt{Q} --- \texttt{finite query set}; \texttt{deterministic residual query universe}
  \par\smallskip\noindent\emph{Definition.} \(\mathcal Q_d=I\times\{-1,+1\}\).
  \item \texttt{m} --- \texttt{regression family}; \(m_{j,b}:\mathbb R\to\mathbb R\); \texttt{nuisance functional}
  \par\smallskip\noindent\emph{Definition.} \(m_{j,b}(x)=\mathbb E_P[E_{j,b}\mid C_{j,b}=x]\) for \((j,b)\in\mathcal Q_d\).
  \item \texttt{r} --- \texttt{residual family}; \texttt{generated population residual}
  \par\smallskip\noindent\emph{Definition.} \(r_{j,b}=E_{j,b}-m_{j,b}(C_{j,b})\).
  \item \texttt{k} --- \texttt{kernel}; \(k:\mathbb R^2\to(0,1]\); \texttt{fixed predictor kernel}
  \par\smallskip\noindent\emph{Definition.} \(k(x,x')=\exp\{-(x-x')^2/(2\sigma_K^2)\}\).
  \item \texttt{ell} --- \texttt{kernel}; \(\ell:\mathbb R^2\to(0,1]\); \texttt{fixed residual kernel}
  \par\smallskip\noindent\emph{Definition.} \(\ell(e,e')=\exp\{-(e-e')^2/(2\sigma_K^2)\}\).
  \item \texttt{muK} --- \texttt{RKHS mean embedding}; \(\mu_R=\mathbb E_R[(k\otimes\ell)((U,V),\cdot)]\); \texttt{dependence-\allowbreak{}functional primitive}
  \par\smallskip\noindent\emph{Definition.} Mean embedding of a law \(R\) into \(\mathcal H_{k\otimes\ell}\).
  \item \texttt{H} --- \texttt{nonnegative functional family}; \(\mathbb R_+^{\mathcal Q_d}\); \texttt{dependence margin}
  \par\smallskip\noindent\emph{Definition.} \(H=(H_{j,b}:(j,b)\in\mathcal Q_d)\) denotes fixed-kernel population residual HSIC.
  \item \texttt{q} --- \texttt{functional family}; \([0,1]^{\mathcal Q_d}\); \texttt{GENE regression coordinate}
  \par\smallskip\noindent\emph{Definition.} \(q=(q_{j,b}:(j,b)\in\mathcal Q_d)\) denotes the scale-normalized explained variance.
  \item \texttt{score} --- \texttt{functional family}; \(\mathbb R^{\mathcal Q_d}\); \texttt{population score}
  \par\smallskip\noindent\emph{Definition.} \(s=(s_{j,b}:(j,b)\in\mathcal Q_d)\) denotes the pairwise population GENE-inspired order scores.
  \item \texttt{g} --- \texttt{vector}; \(\mathbb R^d\); \texttt{navigation margin}
  \par\smallskip\noindent\emph{Definition.} \(g=(g_j:j\in I)\) denotes the population GENE-inspired correcting gains.
  \item \texttt{Gamma} --- \texttt{functional}; \(\Gamma_d:\mathcal M_d\to\mathbb R^{3d}\); \texttt{margin map}
  \par\smallskip\noindent\emph{Definition.} \texttt{Joint map collecting GENE-\allowbreak{}inspired gains and every query-\allowbreak{}indexed residual HSIC value.}
  \item \texttt{PiTheta} --- \texttt{coordinate permutation}; \(\Pi_\theta:\mathbb R^{3d}\to\mathbb R^{3d}\); \texttt{orientation-\allowbreak{}indexed margin-\allowbreak{}coordinate placement}
  \par\smallskip\noindent\emph{Definition.} \(\Pi_\theta((u_j)_{j\in I},(v_{j,+})_{j\in I},(v_{j,-})_{j\in I})=((u_j)_{j\in I},(v_{j,\theta_j})_{j\in I},(v_{j,-\theta_j})_{j\in I})\); thus each ordered block \((H_{j,+},H_{j,-})\) is sent to \((H_{j,\theta_j},H_{j,-\theta_j})\).
  \item \texttt{Fd} --- \texttt{set}; \(\mathbb R^{3d}\); \texttt{focal geometry}
  \par\smallskip\noindent\emph{Definition.} \(\mathfrak F_d\), the feasible image of the joint population margin map.
  \item \texttt{Msep} --- \texttt{model class}; \texttt{minimax parameter space}
  \par\smallskip\noindent\emph{Definition.} \(\mathcal M_d(\lambda)\), the separated product-ANM experiment.
  \item \texttt{a} --- \texttt{vector}; \((0,1)^d\); \texttt{chart coordinate}
  \par\smallskip\noindent\emph{Definition.} \(a=(a_j:j\in I)\), the standardized Gaussian correlation coordinate.
  \item \texttt{t} --- \texttt{vector}; \([0,\infty)^d\); \texttt{chart coordinate}
  \par\smallskip\noindent\emph{Definition.} \(t=(t_j:j\in I)\), the nonlinear Hermite perturbation coordinate.
  \item \texttt{tpath} --- \texttt{scalar-\allowbreak{}coordinate path}; \(t^{[j]}:[0,\infty)\to[0,\infty)^d\); \texttt{dimension-\allowbreak{}correct local chart path}
  \par\smallskip\noindent\emph{Definition.} For \(j\in I\), \(t^{[j]}(u)_j=u\) and \(t^{[j]}(u)_k=0\) for every \(k\ne j\); all coordinates other than the scalar \(j\)-coordinate are fixed along the path.
  \item \texttt{phi} --- \texttt{function}; \(\phi:\mathbb R\to\mathbb R\); \texttt{chart direction}
  \par\smallskip\noindent\emph{Definition.} \(\phi(z)=z^2-1\), the second centered Gaussian Hermite polynomial.
  \item \texttt{Pchart} --- \texttt{law family}; \texttt{least-\allowbreak{}favorable construction}
  \par\smallskip\noindent\emph{Definition.} \(P_{a,t,\theta}\), the oriented Hermite perturbation chart.
  \item \texttt{kappaH} --- \texttt{function}; \(\kappa_H:(0,1)\to(0,\infty)\); \texttt{leading reverse-\allowbreak{}HSIC coefficient}
  \par\smallskip\noindent\emph{Definition.} For \(v\in(0,1)\), fix any \(j\in I\), orientation \(\theta\), and \(a\) with \(a_j=v\), and set \(\kappa_H(v)=\lim_{u\downarrow0}u^{-2}H_{j,-\theta_j}(P_{a,t^{[j]}(u),\theta})\). Along \(t^{[j]}(u)\), every coordinate other than \(j\) is fixed at zero; factorization of \(P_{a,t,\theta}\) makes the displayed coordinate functional independent of the remaining entries of \(a\).
  \item \texttt{rn} --- \texttt{rate}; \texttt{one-\allowbreak{}dimensional regression rate}
  \par\smallskip\noindent\emph{Definition.} \(r_n=n^{-\beta/(2\beta+1)}\).
  \item \texttt{pLP} --- \texttt{integer}; \(\{0,1,\ldots\}\); \texttt{nuisance-\allowbreak{}protocol constant}
  \par\smallskip\noindent\emph{Definition.} \(p_{\mathrm{LP}}=\lfloor\beta\rfloor\), the fixed local-polynomial order.
  \item \texttt{Folds} --- \texttt{ordered partition}; \texttt{cross-\allowbreak{}fitting allocation}
  \par\smallskip\noindent\emph{Definition.} \(\mathcal I_n=(I_1,I_2,I_3,I_4)\) is the deterministic balanced four-fold partition of \(\{1,\ldots,n\}\) by observation index, with fold sizes differing by at most one.
  \item \texttt{hn} --- \texttt{positive scalar}; \texttt{nuisance-\allowbreak{}protocol bandwidth}
  \par\smallskip\noindent\emph{Definition.} \(h_n=(\max\{1,\lfloor n/4\rfloor\})^{-1/(2\beta+1)}\), the deterministic local-polynomial bandwidth.
  \item \texttt{Tn} --- \texttt{positive scalar}; \texttt{nuisance-\allowbreak{}protocol truncation}
  \par\smallskip\noindent\emph{Definition.} \(T_n=\sqrt{8\log(n\vee3)/c_0}\), the deterministic predictor truncation radius.
  \item \texttt{Bn} --- \texttt{positive scalar}; \texttt{nuisance-\allowbreak{}protocol clipping}
  \par\smallskip\noindent\emph{Definition.} \(B_n=L(1+T_n)^2\), the deterministic response-clipping radius.
  \item \texttt{piT} --- \texttt{function}; \(\pi_{T_n}:\mathbb R\to[-T_n,T_n]\); \texttt{tail winsorization map}
  \par\smallskip\noindent\emph{Definition.} \(\pi_{T_n}(z)=\max\{-T_n,\min(z,T_n)\}\).
  \item \texttt{KLP} --- \texttt{kernel}; \(K_{\mathrm{LP}}:\mathbb R\to[0,\infty)\); \texttt{nuisance-\allowbreak{}protocol kernel}
  \par\smallskip\noindent\emph{Definition.} \(K_{\mathrm{LP}}(v)=\tfrac34(1-v^2)\mathbf 1\{|v|\le1\}\), the fixed Epanechnikov local-polynomial kernel.
  \item \texttt{Afit} --- \texttt{estimator class}; \texttt{admissible nuisance-\allowbreak{}estimator class}
  \par\smallskip\noindent\emph{Definition.} \(\mathcal A_n^{\mathrm{LP}}\), the admissible foldwise local-polynomial regression class fixed by Definition \(\mathrm{def:admissible\text{-}local\text{-}polynomial\text{-}class}\).
  \item \texttt{RegEvent} --- \texttt{event family}; \(\mathcal G_n(x;C_m,\widehat m)\); \texttt{uniform nuisance guarantee}
  \par\smallskip\noindent\emph{Definition.} The simultaneous foldwise \(L^4\) regression-accuracy event in Definition \(\mathrm{def:regression\text{-}good\text{-}event}\).
  \item \texttt{OrthHandle} --- \texttt{construction handle}; \texttt{open construction}
  \par\smallskip\noindent\emph{Definition.} \(\mathfrak H_{\mathrm{orth}}(\mathcal A_n^{\mathrm{LP}})\), the second-order residual-embedding orthogonalization route applied only to the admissible four-fold local-polynomial class.
  \item \texttt{Tperp} --- \texttt{target statistic family}; \texttt{open inferential object}
  \par\smallskip\noindent\emph{Definition.} \((\widehat H^{\perp}_{j,b}(\widehat m):(j,b)\in\mathcal Q_d)\), to be derived from \(\mathfrak H_{\mathrm{orth}}(\mathcal A_n^{\mathrm{LP}})\) for \(\widehat m\in\mathcal A_n^{\mathrm{LP}}\).
  \item \texttt{thetahat} --- \texttt{estimator}; \(\{-1,+1\}^I\); \texttt{estimator}
  \par\smallskip\noindent\emph{Definition.} \(\widehat\theta_n\), the simultaneous orientation rule.
  \item \texttt{amin} --- \texttt{scalar}; \((0,1/2)\); \texttt{chart constant}
  \par\smallskip\noindent\emph{Definition.} \(a_-\) is the fixed lower endpoint of the sub-half Hermite chart.
  \item \texttt{amax} --- \texttt{scalar}; \((a_-,1/2)\); \texttt{chart constant}
  \par\smallskip\noindent\emph{Definition.} \(a_+\) is the fixed upper endpoint of the sub-half Hermite chart.
  \item \texttt{Nstar} --- \texttt{integer-\allowbreak{}valued functional}; \texttt{minimax frontier}
  \par\smallskip\noindent\emph{Definition.} \(N^*(d,\alpha,\lambda)\), the smallest sample size allowing worst-case exact-orientation error at most \(\alpha\) over \(\mathcal M_d(\lambda)\).
  \item \texttt{tstar} --- \texttt{scalar}; \((0,\infty)\); \texttt{chart constant}
  \par\smallskip\noindent\emph{Definition.} \(t_\star\) is the fixed positive radius of the calibrated sub-half chart.
  \item \texttt{lambdaH} --- \texttt{scalar}; \((0,1]\); \texttt{chart signal radius}
  \par\smallskip\noindent\emph{Definition.} \(\lambda_H\le\lambda_0\) is the local signal radius of the sub-half Hermite chart.
  \item \texttt{Mchart} --- \texttt{model class}; \texttt{matched minimax parameter space}
  \par\smallskip\noindent\emph{Definition.} \(\mathcal M_d^{\mathrm H}(\lambda)\), the separated sub-half Hermite-chart experiment.
  \item \texttt{psi} --- \texttt{function}; \(\psi:\mathbb R^2\to\mathbb R\); \texttt{orientation statistic}
  \par\smallskip\noindent\emph{Definition.} \(\psi(x,y)=(x^2-1)y-(y^2-1)x\), the antisymmetric observable chart moment.
  \item \texttt{Bset} --- \texttt{set-\allowbreak{}valued decision family}; \(2^{\{-1,+1\}}\); \texttt{confidence-\allowbreak{}set coordinates}
  \par\smallskip\noindent\emph{Definition.} \((B_j:j\in I)\), the pairwise sign-acceptance sets formed by simultaneous inversion of the chart moment intervals.
  \item \texttt{Chat} --- \texttt{random set}; \(2^{\{-1,+1\}^I}\); \texttt{uncertainty object}
  \par\smallskip\noindent\emph{Definition.} \(\widehat{\mathcal C}_n\), the simultaneous orientation confidence set.
  \item \texttt{Vpsi} --- \texttt{positive scalar}; \texttt{concentration constant}
  \par\smallskip\noindent\emph{Definition.} \(V_\psi\) is the explicit uniform variance envelope in Definition \(\mathrm{def:moment\text{-}variance\text{-}constant}\).
  \item \texttt{Nchart} --- \texttt{integer-\allowbreak{}valued functional}; \texttt{chart minimax frontier}
  \par\smallskip\noindent\emph{Definition.} \(N_{\mathrm H}^*(d,\alpha,\lambda)\), the all-procedure exact-orientation sample complexity on \(\mathcal M_d^{\mathrm H}(\lambda)\).
\end{itemize}

\section{Assumptions}

\begin{assumption}[ass:pair-product]\label{ass:pair-product}
\[P=\bigotimes_{j\in I}P_j.\]
\par\smallskip\noindent\emph{Standard} (independent product experiment, \cite{Tsybakov2009}).
\end{assumption}

\begin{assumption}[ass:anm-equations]\label{ass:anm-equations}
\[E_j=f_j(C_j)+\varepsilon_j\quad\text{for every }j\in I.\]
\par\smallskip\noindent\emph{Standard} (continuous additive-noise structural equations, \cite{PetersMJS2014}).
\end{assumption}

\begin{assumption}[ass:noise-exogeneity]\label{ass:noise-exogeneity}
\[\varepsilon_j\perp C_j\quad\text{for every }j\in I.\]
\par\smallskip\noindent\emph{Standard} (ANM noise independence, \cite{PetersMJS2014}).
\end{assumption}

\begin{assumption}[ass:restricted-anm]\label{ass:restricted-anm}
\[(f_j,p^C_j,p^\varepsilon_j)\text{ satisfies Peters et al.'s Condition 19 for every }j\in I.\]
\par\smallskip\noindent\emph{Standard} (Peters restricted-ANM Condition 19, \cite{PetersMJS2014}).
\end{assumption}

\begin{assumption}[ass:holder-envelope]\label{ass:holder-envelope}
\[f_j,\log p^C_j,\log p^\varepsilon_j\in\mathcal H_{(2)}^\beta(L)\quad\text{for every }j\in I.\]
\par\smallskip\noindent\emph{Novel.} The fixed weighted Holder envelope describes smooth mechanisms and smooth log densities on unbounded supports in one common scale. It contains Gaussian causes and noises together with bounded smooth log-density perturbations and polynomially growing nonlinear mechanisms, covering standardized physical measurement pairs whose response curve is smooth and whose tails remain Gaussian-like.
\end{assumption}

\begin{assumption}[ass:gaussian-tail-envelope]\label{ass:gaussian-tail-envelope}
\[c_0e^{-C_0z^2}\le p_j(z)\le C_0e^{-c_0z^2}\quad\text{for every }p_j\in\{p^C_j,p^\varepsilon_j\},\ z\in\mathbb R,\ j\in I.\]
\par\smallskip\noindent\emph{Novel.} The two-sided density envelope retains full support while excluding tail sequences that become arbitrarily sparse. It is satisfied by standardized Gaussian causes and noises and by uniformly bounded smooth exponential tilts of them, a concrete domain for continuously measured pairs with approximately Gaussian latent shocks.
\end{assumption}

\begin{assumption}[ass:cause-unit-variance]\label{ass:cause-unit-variance}
\[
\operatorname{Var}_P(C_j)=1\quad\text{for every }j\in I.
\]
\par\smallskip\noindent\emph{Standard} (cause-coordinate scale normalization for the GENE-inspired score, \cite{pmlr-v267-li25bx}).
\end{assumption}

\begin{assumption}[ass:effect-unit-variance]\label{ass:effect-unit-variance}
\[
\operatorname{Var}_P(E_j)=1\quad\text{for every }j\in I.
\]
\par\smallskip\noindent\emph{Standard} (effect-coordinate scale normalization for the GENE-inspired score, \cite{pmlr-v267-li25bx}).
\end{assumption}

\begin{assumption}[ass:fixed-predictor-kernel]\label{ass:fixed-predictor-kernel}
\[k(x,x')=e^{-(x-x')^2/(2\sigma_K^2)}.\]
\par\smallskip\noindent\emph{Standard} (bounded characteristic Gaussian predictor kernel, \cite{GrettonBSS05}).
\end{assumption}

\begin{assumption}[ass:fixed-residual-kernel]\label{ass:fixed-residual-kernel}
\[\ell(e,e')=e^{-(e-e')^2/(2\sigma_K^2)}.\]
\par\smallskip\noindent\emph{Standard} (bounded characteristic Gaussian residual kernel, \cite{GrettonBSS05}).
\end{assumption}

\begin{assumption}[ass:navigation-window]\label{ass:navigation-window}
\[
\underline g\le g_j\le\overline g\quad\text{for every }j\in I.
\]
\par\smallskip\noindent\emph{Novel.} The navigation window selects supplied pairs with stable, nondegenerate GENE-inspired correcting gains. Moderate-signal systems in which the true mechanism explains an interior fraction of effect variance satisfy the restriction, including the oriented Hermite chart over the calibrated compact sub-half correlation interval.
\end{assumption}

\begin{assumption}[ass:residual-hsic-window]\label{ass:residual-hsic-window}
\[\lambda\le H_{j,-\theta_j}\le A_H\lambda\quad\text{for every }j\in I.\]
\par\smallskip\noindent\emph{Novel.} The residual-HSIC window defines a local batch of pairs whose wrong-direction departures from independence have comparable strength within a fixed factor. Such batches arise when scientific modules are screened to a common detection scale, and the Hermite chart realizes every coordinate of this window while preserving positive densities and restricted-ANM identification.
\end{assumption}

\section{Definitions}

\begin{definition}[def:base-model]\label{def:base-model}
\par\smallskip\noindent\emph{Defined object.} \(\mathcal M_d(\beta,L,c_0,C_0)\)
\par\smallskip\noindent\emph{Construction.} \[\mathcal M_d(\beta,L,c_0,C_0)=\{(P,\theta,f,p^C,p^\varepsilon):d\ge1\text{ and the eight member properties listed below hold}\}.\]
\par\smallskip\noindent\emph{Member properties.} \texttt{ass:\allowbreak{}pair-\allowbreak{}product}; \texttt{ass:\allowbreak{}anm-\allowbreak{}equations}; \texttt{ass:\allowbreak{}noise-\allowbreak{}exogeneity}; \texttt{ass:\allowbreak{}restricted-\allowbreak{}anm}; \texttt{ass:\allowbreak{}holder-\allowbreak{}envelope}; \texttt{ass:\allowbreak{}gaussian-\allowbreak{}tail-\allowbreak{}envelope}; \texttt{ass:\allowbreak{}cause-\allowbreak{}unit-\allowbreak{}variance}; \texttt{ass:\allowbreak{}effect-\allowbreak{}unit-\allowbreak{}variance}.
\end{definition}

\begin{definition}[def:query-universe]\label{def:query-universe}
\par\smallskip\noindent\emph{Defined object.} \texttt{deterministic pairwise query universe}
\par\smallskip\noindent\emph{Construction.} \[\mathcal Q_d=\{(j,b):j\in I,\ b\in\{-1,+1\}\}.\]
\end{definition}

\begin{definition}[def:population-regressions]\label{def:population-regressions}
\par\smallskip\noindent\emph{Defined object.} \texttt{direction-\allowbreak{}indexed conditional means}
\par\smallskip\noindent\emph{Construction.} \[m_{j,b}(x)=\mathbb E_P[E_{j,b}\mid C_{j,b}=x],\qquad (j,b)\in\mathcal Q_d.\]
\end{definition}

\begin{definition}[def:population-residuals]\label{def:population-residuals}
\par\smallskip\noindent\emph{Defined object.} \texttt{direction-\allowbreak{}indexed population residuals}
\par\smallskip\noindent\emph{Construction.} \[r_{j,b}=E_{j,b}-m_{j,b}(C_{j,b}),\qquad (j,b)\in\mathcal Q_d.\]
\par\smallskip\noindent\emph{Inputs.} \texttt{def:\allowbreak{}population-\allowbreak{}regressions}.
\end{definition}

\begin{definition}[def:hsic]\label{def:hsic}
\par\smallskip\noindent\emph{Defined object.} \texttt{fixed-\allowbreak{}kernel residual HSIC}
\par\smallskip\noindent\emph{Construction.} \[H_{j,b}=\bigl\|\mu_{P_{C_{j,b},r_{j,b}}}-\mu_{P_{C_{j,b}}\otimes P_{r_{j,b}}}\bigr\|^2_{\mathcal H_{k\otimes\ell}},\qquad (j,b)\in\mathcal Q_d.\]
\par\smallskip\noindent\emph{Inputs.} \texttt{def:\allowbreak{}population-\allowbreak{}residuals}; \texttt{ass:\allowbreak{}fixed-\allowbreak{}predictor-\allowbreak{}kernel}; \texttt{ass:\allowbreak{}fixed-\allowbreak{}residual-\allowbreak{}kernel}.
\end{definition}

\begin{definition}[def:gene-gain]\label{def:gene-gain}
\par\smallskip\noindent\emph{Defined object.} \texttt{GENE-\allowbreak{}inspired population correcting gain}
\par\smallskip\noindent\emph{Construction.} \[
q_{j,b}=\frac{\operatorname{Var}_P\{m_{j,b}(C_{j,b})\}}{\operatorname{Var}_P(E_{j,b})},\qquad
s_{j,b}=q_{j,b}\left(1-\frac12\mathbf 1\{H_{j,b}>0\}\right),\qquad
g_j=s_{j,\theta_j}-s_{j,-\theta_j}.
\]
This is a GENE-inspired population score: its independence coordinate is the oracle indicator \(\mathbf 1\{H_{j,b}>0\}\) for fixed characteristic Gaussian-kernel HSIC.  It is not the published GENE score, which uses a discretized chi-square independence-test indicator together with neural-network \(R^2\).
\par\smallskip\noindent\emph{Inputs.} \texttt{def:\allowbreak{}population-\allowbreak{}regressions}; \texttt{def:\allowbreak{}hsic}; \texttt{ass:\allowbreak{}cause-\allowbreak{}unit-\allowbreak{}variance}; \texttt{ass:\allowbreak{}effect-\allowbreak{}unit-\allowbreak{}variance}.
\end{definition}

\begin{definition}[def:margin-map]\label{def:margin-map}
\par\smallskip\noindent\emph{Defined object.} \texttt{joint GENE-\allowbreak{}inspired-\allowbreak{}-\allowbreak{}HSIC margin map}
\par\smallskip\noindent\emph{Construction.} \[
\Gamma_d(P,\theta)=\bigl((g_j)_{j\in I},(H_{j,+})_{j\in I},(H_{j,-})_{j\in I}\bigr).
\]
\par\smallskip\noindent\emph{Inputs.} \texttt{def:\allowbreak{}hsic}; \texttt{def:\allowbreak{}gene-\allowbreak{}gain}.
\end{definition}

\begin{definition}[def:orientation-permutation]\label{def:orientation-permutation}
\par\smallskip\noindent\emph{Defined object.} \texttt{orientation-\allowbreak{}indexed HSIC-\allowbreak{}block permutation}
\par\smallskip\noindent\emph{Construction.} \[\Pi_\theta\bigl((u_j)_{j\in I},(v_{j,+})_{j\in I},(v_{j,-})_{j\in I}\bigr)=\bigl((u_j)_{j\in I},(v_{j,\theta_j})_{j\in I},(v_{j,-\theta_j})_{j\in I}\bigr).\]
\par\smallskip\noindent\emph{Inputs.} \texttt{def:\allowbreak{}margin-\allowbreak{}map}.
\end{definition}

\begin{definition}[def:feasible-set]\label{def:feasible-set}
\par\smallskip\noindent\emph{Defined object.} \(\mathfrak F_d\)
\par\smallskip\noindent\emph{Construction.} \[\mathfrak F_d=\{\Gamma_d(P,\theta):(P,\theta,f,p^C,p^\varepsilon)\in\mathcal M_d(\beta,L,c_0,C_0)\}.\]
\par\smallskip\noindent\emph{Inputs.} \texttt{def:\allowbreak{}margin-\allowbreak{}map}; \texttt{def:\allowbreak{}base-\allowbreak{}model}; \texttt{ass:\allowbreak{}fixed-\allowbreak{}predictor-\allowbreak{}kernel}; \texttt{ass:\allowbreak{}fixed-\allowbreak{}residual-\allowbreak{}kernel}.
\end{definition}

\begin{definition}[def:separated-model]\label{def:separated-model}
\par\smallskip\noindent\emph{Defined object.} \(\mathcal M_d(\lambda)\)
\par\smallskip\noindent\emph{Construction.} \[\mathcal M_d(\lambda)=\{(P,\theta,f,p^C,p^\varepsilon)\in\mathcal M_d(\beta,L,c_0,C_0):\underline g\le g_j\le\overline g,\ \lambda\le H_{j,-\theta_j}\le A_H\lambda\text{ for every }j\in I\}.\]
\par\smallskip\noindent\emph{Member properties.} \texttt{ass:\allowbreak{}navigation-\allowbreak{}window}; \texttt{ass:\allowbreak{}residual-\allowbreak{}hsic-\allowbreak{}window}.
\end{definition}

\begin{definition}[def:scalar-coordinate-path]\label{def:scalar-coordinate-path}
\par\smallskip\noindent\emph{Defined object.} \texttt{one-\allowbreak{}coordinate Hermite path}
\par\smallskip\noindent\emph{Construction.} \[t^{[j]}(u)_k=u\mathbf 1\{k=j\},\qquad u\ge0,\quad j,k\in I.\] Thus coordinate \(j\) varies with the scalar \(u\), while every other coordinate remains fixed at zero.
\end{definition}

\begin{definition}[def:hermite-chart]\label{def:hermite-chart}
\par\smallskip\noindent\emph{Defined object.} \texttt{oriented Hermite perturbation chart}
\par\smallskip\noindent\emph{Construction.} For \(a_j\in(0,1)\) and \(0\le t_j<\sqrt{(1-a_j^2)/2}\), draw independent \(Z_j,\xi_j\sim\mathcal N(0,1)\), set \(F_{a_j,t_j}(z)=a_jz+t_j\phi(z)\) and \(s_{a_j,t_j}=\sqrt{1-a_j^2-2t_j^2}\), and define \((X_j,Y_j)=(Z_j,F_{a_j,t_j}(Z_j)+s_{a_j,t_j}\xi_j)\) when \(\theta_j=+1\) and the coordinate-swapped pair when \(\theta_j=-1\); \(P_{a,t,\theta}=\bigotimes_{j\in I}\mathcal L(X_j,Y_j)\). Boundary values with \(t_k=0\) are included in the chart domain to define scalar-coordinate limits.
\par\smallskip\noindent\emph{Inputs.} \texttt{def:\allowbreak{}query-\allowbreak{}universe}; \texttt{def:\allowbreak{}scalar-\allowbreak{}coordinate-\allowbreak{}path}.
\end{definition}

\begin{definition}[def:balanced-four-fold-allocation]\label{def:balanced-four-fold-allocation}
\par\smallskip\noindent\emph{Defined object.} \texttt{balanced cyclic four-\allowbreak{}fold allocation}
\par\smallskip\noindent\emph{Construction.} Let \(\mathcal I_{n,r}=\{i\in\{1,\ldots,n\}:i\equiv r\pmod 4\}\), with residues relabelled as \(r\in\{1,2,3,4\}\). For rotation \(r\), folds \((\mathcal I_{n,r},\mathcal I_{n,r+1},\mathcal I_{n,r+2},\mathcal I_{n,r+3})\), with indices read modulo four, are assigned respectively to regression fitting, derivative-projection fitting, the first corrected-feature mean, and the second corrected-feature mean; the four rotation-specific statistics are averaged.
\end{definition}

\begin{definition}[def:admissible-local-polynomial-class]\label{def:admissible-local-polynomial-class}
\par\smallskip\noindent\emph{Defined object.} \(\mathcal A_n^{\mathrm{LP}}\)
\par\smallskip\noindent\emph{Construction.} \[\mathcal A_n^{\mathrm{LP}}=\left\{\widehat m=(\widehat m^{(r)}_{j,b}:r\in\{1,2,3,4\},(j,b)\in\mathcal Q_d):\widehat m\text{ is produced by the following deterministic protocol}\right\}.\] On regression fold \(\mathcal I_{n,r}\), for \(x\in\mathbb R\), set \(z=\pi_{T_n}(x)\) and let \((\widehat a_0,\ldots,\widehat a_{p_{\mathrm{LP}}})\) be the unique minimizer of \[\sum_{i\in\mathcal I_{n,r}:\,|C_{j,b}(W_i)|\le T_n}K_{z,n}\!\left(\frac{C_{j,b}(W_i)-z}{h_n}\right)\left[E_{j,b}(W_i)-\sum_{q=0}^{p_{\mathrm{LP}}}a_q(C_{j,b}(W_i)-z)^q\right]^2+n^{-2}\sum_{q=0}^{p_{\mathrm{LP}}}a_q^2.\] Here \(K_{z,n}(v)=K_{\mathrm{LP}}(v)\) for \(|z|\le T_n-h_n\), \(K_{z,n}(v)=2K_{\mathrm{LP}}(v)\mathbf 1\{v\ge0\}\) for \(z<-T_n+h_n\), and \(K_{z,n}(v)=2K_{\mathrm{LP}}(v)\mathbf 1\{v\le0\}\) for \(z>T_n-h_n\). Set \(\widehat m^{(r)}_{j,b}(x)=\operatorname{clip}_{[-B_n,B_n]}(\widehat a_0)\). The ridge makes the minimizer unique, so the displayed set is a nonempty singleton collection of all four rotation-specific fits; no other nuisance fit is admissible.
\par\smallskip\noindent\emph{Inputs.} \texttt{def:\allowbreak{}balanced-\allowbreak{}four-\allowbreak{}fold-\allowbreak{}allocation}; \texttt{def:\allowbreak{}query-\allowbreak{}universe}; \texttt{def:\allowbreak{}population-\allowbreak{}regressions}.
\end{definition}

\begin{definition}[def:regression-good-event]\label{def:regression-good-event}
\par\smallskip\noindent\emph{Defined object.} \texttt{simultaneous admissible-\allowbreak{}regression event}
\par\smallskip\noindent\emph{Construction.} For \(x\ge1\), \(u=x+\log(2d)\), \(C_m>0\), and \(\widehat m\in\mathcal A_n^{\mathrm{LP}}\), define \[\mathcal G_n(x;C_m,\widehat m)=\left\{\max_{r\in\{1,2,3,4\}}\max_{(j,b)\in\mathcal Q_d}\|\widehat m^{(r)}_{j,b}-m_{j,b}\|_{L^4(P_{C_{j,b}})}\le C_m\left(r_n+\sqrt{\frac{u}{n}}\right)\right\}.\]
\par\smallskip\noindent\emph{Inputs.} \texttt{def:\allowbreak{}admissible-\allowbreak{}local-\allowbreak{}polynomial-\allowbreak{}class}; \texttt{def:\allowbreak{}population-\allowbreak{}regressions}.
\end{definition}

\begin{definition}[def:orthogonalization-handle]\label{def:orthogonalization-handle}
\par\smallskip\noindent\emph{Defined object.} \texttt{second-\allowbreak{}order residual-\allowbreak{}embedding orthogonalization handle}
\par\smallskip\noindent\emph{Construction.} For each \(\widehat m\in\mathcal A_n^{\mathrm{LP}}\) and each cyclic rotation of \(\mathcal I_n\), form generated residuals from the designated regression fold, project the first Gateaux derivative of the centered Gaussian-kernel residual feature map on the conditional-mean tangent using only the designated derivative-projection fold, form two corrected feature means on the two remaining folds, cross-multiply them, remove the estimated second-order diagonal term using the derivative-projection fold, studentize with held-out conditional variances, and average the four rotation-specific statistics. The construction must be measurable in the four allocated folds and is evaluated on \(\mathcal G_n(x;C_m,\widehat m)\).
\par\smallskip\noindent\emph{Inputs.} \texttt{def:\allowbreak{}population-\allowbreak{}regressions}; \texttt{def:\allowbreak{}population-\allowbreak{}residuals}; \texttt{def:\allowbreak{}hsic}; \texttt{def:\allowbreak{}balanced-\allowbreak{}four-\allowbreak{}fold-\allowbreak{}allocation}; \texttt{def:\allowbreak{}admissible-\allowbreak{}local-\allowbreak{}polynomial-\allowbreak{}class}; \texttt{def:\allowbreak{}regression-\allowbreak{}good-\allowbreak{}event}.
\end{definition}

\begin{definition}[def:minimax-frontier]\label{def:minimax-frontier}
\par\smallskip\noindent\emph{Defined object.} \texttt{exact-\allowbreak{}orientation minimax sample complexity}
\par\smallskip\noindent\emph{Construction.} \[N^*(d,\alpha,\lambda)=\inf\left\{n:\inf_{\widetilde\theta}\sup_{(P,\theta)\in\mathcal M_d(\lambda)}P^{\otimes n}(\widetilde\theta\ne\theta)\le\alpha\right\}.\]
\end{definition}

\begin{definition}[def:subhalf-hermite-class]\label{def:subhalf-hermite-class}
\par\smallskip\noindent\emph{Defined object.} \texttt{separated sub-\allowbreak{}half Hermite-\allowbreak{}chart experiment}
\par\smallskip\noindent\emph{Construction.} Fix constants \(0<a_-<a_+<1/2\), \(t_\star>0\), and \(0<\lambda_H\le\lambda_0\).  For \(0<\lambda\le\lambda_H\), define
\[
\mathcal M_d^{\mathrm H}(\lambda)
=\left\{(P_{a,t,\theta},\theta):
\begin{array}{l}
\theta\in\{-1,+1\}^{I},\ a\in[a_-,a_+]^d,\ 0<t_j\le t_\star\ (j\in I),\\
(P_{a,t,\theta},\theta)\text{ is the law--orientation projection of a member of }\mathcal M_d(\lambda)
\end{array}
\right\}.
\]
Thus \(\mathcal M_d^{\mathrm H}(\lambda)\) is an explicitly parameterized subclass of the separated positive-density restricted-ANM experiment; membership still requires every law property in \(\mathcal M_d(\lambda)\).
\par\smallskip\noindent\emph{Inputs.} \texttt{def:\allowbreak{}hermite-\allowbreak{}chart}; \texttt{def:\allowbreak{}separated-\allowbreak{}model}.
\end{definition}

\begin{definition}[def:chart-moment-observable]\label{def:chart-moment-observable}
\par\smallskip\noindent\emph{Defined object.} \texttt{antisymmetric chart moment}
\par\smallskip\noindent\emph{Construction.} For one labelled pair define the antisymmetric observable
\[
\psi(x,y)=(x^2-1)y-(y^2-1)x.
\]
For observation \(W_i\) and pair \(j\), write \(\psi_{ij}=\psi(X_j(W_i),Y_j(W_i))\).
\end{definition}

\begin{definition}[def:moment-variance-constant]\label{def:moment-variance-constant}
\par\smallskip\noindent\emph{Defined object.} \texttt{computable chart-\allowbreak{}moment variance envelope}
\par\smallskip\noindent\emph{Construction.} Let \(Z\sim\mathcal N(0,1)\), put
\[
M_8=3^7\left(2\mathbb E|Z|^8+\mathbb E|Z^2-1|^8\right),
\]
and define the computable constant
\[
V_\psi=4\left\{(\mathbb E|Z|^8)^{1/2}M_8^{1/4}+2+M_8^{1/2}(\mathbb E|Z|^4)^{1/2}\right\}.
\]
All Gaussian moments in this display are explicit integers or finite sums of explicit integers.
\par\smallskip\noindent\emph{Inputs.} \texttt{def:\allowbreak{}chart-\allowbreak{}moment-\allowbreak{}observable}.
\end{definition}

\begin{definition}[def:chart-moment-rule]\label{def:chart-moment-rule}
\par\smallskip\noindent\emph{Defined object.} \texttt{median-\allowbreak{}of-\allowbreak{}means chart orientation rule and confidence set}
\par\smallskip\noindent\emph{Construction.} Let
\[
K=\left\lceil 8\log\frac{2d}{\alpha}\right\rceil,
\qquad q=\left\lfloor\frac nK\right\rfloor,
\qquad \rho_{n,d,\alpha}=\sqrt{\frac{8V_\psi}{q}},
\]
and suppose \(n\ge2K\).  Partition the first \(Kq\) observation indices deterministically into \(K\) consecutive blocks \(A_1,\ldots,A_K\) of size \(q\).  For each \(j\in I\), set
\[
\widehat\mu_j=\operatorname{Med}_{1\le k\le K}
\left\{q^{-1}\sum_{i\in A_k}\psi_{ij}\right\},
\]
using the lower median when \(K\) is even, and form
\[
B_j=\{+1:\widehat\mu_j+\rho_{n,d,\alpha}\ge0\}
\cup\{-1:\widehat\mu_j-\rho_{n,d,\alpha}\le0\}.
\]
Set \(\widehat{\mathcal C}_n^{\mathrm H}=\prod_{j\in I}B_j\).  Let \(\widehat\theta_n^{\mathrm H}\) be the unique element of this product when it is a singleton and a fixed measurable element otherwise.  For \(n<2K\), define \(B_j=\{-1,+1\}\) and use the same fixed fallback.  This map uses only the \(d\) observable moment queries \(\psi_{ij}\).
\par\smallskip\noindent\emph{Inputs.} \texttt{def:\allowbreak{}chart-\allowbreak{}moment-\allowbreak{}observable}; \texttt{def:\allowbreak{}moment-\allowbreak{}variance-\allowbreak{}constant}.
\end{definition}

\begin{definition}[def:chart-minimax-frontier]\label{def:chart-minimax-frontier}
\par\smallskip\noindent\emph{Defined object.} \texttt{sub-\allowbreak{}half Hermite-\allowbreak{}chart exact-\allowbreak{}orientation sample complexity}
\par\smallskip\noindent\emph{Construction.} \[
N_{\mathrm H}^*(d,\alpha,\lambda)=
\inf\left\{n:\inf_{\widetilde\theta}
\sup_{(P,\theta)\in\mathcal M_d^{\mathrm H}(\lambda)}
P^{\otimes n}(\widetilde\theta\ne\theta)\le\alpha\right\},
\]
where the inner infimum ranges over every measurable map from \((\mathbb R^{2d})^n\) to \(\{-1,+1\}^{I}\).
\par\smallskip\noindent\emph{Inputs.} \texttt{def:\allowbreak{}subhalf-\allowbreak{}hermite-\allowbreak{}class}.
\end{definition}

\section{Main results}

\begin{lemma}[lem:population-direction-margin]\label{lem:population-direction-margin}
For every \((P,\theta,f,p^C,p^\varepsilon)\in\mathcal M_d(\beta,L,c_0,C_0)\), \(H_{j,\theta_j}=0\) for all \(j\in I\). If the mechanism is non-affine, Peters et al.'s restricted-ANM condition rules out the reverse additive-noise representation, so characteristic Gaussian kernels give \(H_{j,-\theta_j}>0\), and \(\theta_j\) is the unique population direction with zero residual HSIC.
\end{lemma}
\begin{proof}
Fix \(j\in I\).  By Assumptions \(\mathrm{ass:anm\text{-}equations}\) and \(\mathrm{ass:noise\text{-}exogeneity}\),
\[
m_{j,\theta_j}(C_j)=f_j(C_j)+\mathbb E_P\varepsilon_j,
\qquad
r_{j,\theta_j}=\varepsilon_j-\mathbb E_P\varepsilon_j\perp C_j.
\]
Thus the joint law of \((C_j,r_{j,\theta_j})\) equals the product of its marginals.  Definition \(\mathrm{def:hsic}\) then gives \(H_{j,\theta_j}=0\).

For completeness, the Gaussian-kernel implication used in the reverse direction follows directly from Fourier analysis.  If \(\nu=R_{UV}-R_U\otimes R_V\), then Fubini's theorem, applicable because the kernels are bounded, gives
\[
\left\|\mu_R-\mu_{R_U\otimes R_V}\right\|_{\mathcal H_{k\otimes\ell}}^2
=c_{\sigma_K}\int_{\mathbb R^2}|\widehat\nu(\omega_1,\omega_2)|^2
e^{-\sigma_K^2(\omega_1^2+\omega_2^2)/2}\,d\omega_1d\omega_2.
\]
The weight is strictly positive.  Hence a zero norm implies \(\widehat\nu=0\) almost everywhere and, by continuity, everywhere.  Convolving \(\nu\) with a centered Gaussian of variance \(\eta^2\) gives a signed density whose Fourier transform is zero, so Fourier inversion makes every such convolution zero.  Letting \(\eta\downarrow0\) and testing against bounded continuous functions gives \(\nu=0\).  Consequently, under Assumptions \(\mathrm{ass:fixed\text{-}predictor\text{-}kernel}\) and \(\mathrm{ass:fixed\text{-}residual\text{-}kernel}\), \(H_{j,b}=0\) if and only if \(C_{j,b}\perp r_{j,b}\).

Suppose now that \(f_j\) is non-affine and, contrary to the claim, \(H_{j,-\theta_j}=0\).  The preceding equivalence and Definitions \(\mathrm{def:population\text{-}regressions}\) and \(\mathrm{def:population\text{-}residuals}\) then produce a reverse additive-noise equation
\[
C_j=m_{j,-\theta_j}(E_j)+r_{j,-\theta_j},
\qquad r_{j,-\theta_j}\perp E_j.
\]
Write \(\xi=\log p_j^C\) and \(\nu=\log p_j^\varepsilon\).  Differentiating the logarithm of the positive joint density in the forward and reverse factorizations and eliminating the reverse log densities yields, wherever \(\nu''(y-f_j(x))f_j'(x)\ne0\),
\[
\xi'''=\xi''\left(-\frac{\nu'''f_j'}{\nu''}+\frac{f_j''}{f_j'}\right)
-2\nu''f_j''f_j'+\nu'f_j'''
+\frac{\nu'\nu'''f_j''f_j'}{\nu''}
-\frac{\nu'(f_j'')^2}{f_j'}.
\]
This is exactly the differential equation excluded by Peters et al.'s Condition 19, which is imposed pairwise by Assumption \(\mathrm{ass:restricted\text{-}anm}\).  Thus the reverse representation is impossible.  Since \(H_{j,-\theta_j}\) is a squared Hilbert-space norm, it is nonnegative, so it must be strictly positive.  Therefore \(\theta_j\) is the unique zero-HSIC direction.
\end{proof}
\par\smallskip\noindent \textit{Justification.} This separates qualitative direction identification from the quantitative margins used by the statistical frontier. \textit{Gap.} Peters et al. (2014, Definition 27 and Theorem 28) establish uniqueness, while Mooij et al. (2016, Appendix A) use this residual asymmetry for pointwise consistency; neither result quantifies the joint GENE--HSIC image. \textit{Consumer.} Every later orientation rule can use the deterministic \(2d\)-query universe without enumerating causal orders.

\begin{theorem}[thm:attainable-margin-rectangle]\label{thm:attainable-margin-rectangle}
There exist fixed \(0<a_-<a_+<1\), \(t_0>0\), \(0<\underline g<\overline g<1\), \(A_H>1\), \(\lambda_0>0\), and an envelope \((\beta,L,c_0,C_0)\) such that, for every fixed \(\theta\in\{-1,+1\}^I\), each chart law \(P_{a,t,\theta}\) with \(a\in[a_-,a_+]^d\) and \(0<t_j\le t_0\) belongs to \(\mathcal M_d(\beta,L,c_0,C_0)\). For each \(j\), the scalar-coordinate path \(u\mapsto P_{a,t^{[j]}(u),\theta}\), with every coordinate \(k\ne j\) fixed at zero, satisfies uniformly in \(a_j\) \[H_{j,-\theta_j}(P_{a,t^{[j]}(u),\theta})=\kappa_H(a_j)u^2+O(u^3),\qquad g_j(P_{a,t^{[j]}(u),\theta},\theta)=\frac{a_j^2}{2}+O(u^2),\] where \(\kappa_H(a_j)=\lim_{u\downarrow0}u^{-2}H_{j,-\theta_j}(P_{a,t^{[j]}(u),\theta})\), \(\inf_{v\in[a_-,a_+]}\kappa_H(v)>0\), and the one-sided Jacobian of the full positive chart coordinate \((a_j,t_j^2)\mapsto(g_j,H_{j,-\theta_j})\) has determinant bounded away from zero. Consequently, for every \(0<\lambda\le\lambda_0\), \[[\underline g,\overline g]^d\times\{0\}^d\times[\lambda,A_H\lambda]^d\subseteq\left\{\Pi_\theta\bigl(\Gamma_d(P_{a,t,\theta},\theta)\bigr):a\in[a_-,a_+]^d,\ 0<t_j\le t_0\text{ for every }j\in I\right\}\subseteq\Pi_\theta(\mathfrak F_d).\]
\end{theorem}
\begin{proof}
It suffices first to work with one pair in orientation \(+1\).  Fix
\[
A=[a_-,a_+]=[1/8,1/4]\subset(0,1/2)\subset(0,1/\sqrt 3)
\]
and, for \(a\in A\), let
\[
Z,\xi\stackrel{\mathrm{ind}}{\sim}N(0,1),\qquad
Y_t=aZ+t(Z^2-1)+s_{a,t}\xi,\qquad
s_{a,t}^2=1-a^2-2t^2.
\]
Choose \(t_0>0\) so small that
\[
\underline s^2:=\inf_{a\in A,\,0\le t\le t_0}s_{a,t}^2>0.
\]
The cause \(Z\) and structural noise \(s_{a,t}\xi\) are independent, their densities are strictly positive, and the structural equation is the one in Assumptions \(\mathrm{ass:anm\text{-}equations}\) and \(\mathrm{ass:noise\text{-}exogeneity}\), with
\[
f_{a,t}(z)=az+t(z^2-1).
\]
Moreover,
\[
\operatorname{Var}(Z)=1,\qquad
\operatorname{Var}(Y_t)=a^2+2t^2+s_{a,t}^2=1,
\]
so the two scale-normalization assumptions hold.

We next verify the fixed envelope and the restricted-ANM property rather than treating them as consequences of the chart definition.  Take \(\beta=3\).  The function \(f_{a,t}\), the standard-normal log density, and the \(N(0,s_{a,t}^2)\) log density are polynomials of degree at most two.  Their derivatives through order three, after multiplication by the weight \((1+|z|)^{-2}\) appearing in \(\mathcal H_{(2)}^3(L)\), are bounded uniformly over \(a\in A\) and \(0\le t\le t_0\), and their order-three H\"older seminorm is zero.  Hence one finite \(L>1\) makes Assumption \(\mathrm{ass:holder\text{-}envelope}\) hold throughout the chart.  Since \(s_{a,t}^2\in[\underline s^2,1]\), the Gaussian density formula permits constants \(c_0\in(0,1)\) and \(C_0>1\), independent of \(a,t\), such that
\[
c_0e^{-C_0z^2}\le (2\pi v)^{-1/2}e^{-z^2/(2v)}\le C_0e^{-c_0z^2},\qquad
v\in\{1,s_{a,t}^2\},\ z\in\mathbb R.
\]
Thus Assumption \(\mathrm{ass:gaussian\text{-}tail\text{-}envelope}\) also holds.

For completeness, write \(\xi_C=\log p^C\), \(\nu=\log p^\varepsilon\), and \(e=y-f_{a,t}(z)\).  The differential identity excluded by Peters et al.'s Condition 19, at points where \(f'_{a,t}\nu''\ne0\), is
\[
\xi_C'''=\xi_C''\left(-\frac{\nu'''f'_{a,t}}{\nu''}+\frac{f''_{a,t}}{f'_{a,t}}\right)
-2\nu''f''_{a,t}f'_{a,t}+\nu'f'''_{a,t}
+\frac{\nu'\nu'''f''_{a,t}f'_{a,t}}{\nu''}
-\frac{\nu'(f''_{a,t})^2}{f'_{a,t}}.
\]
Here
\[
\xi_C''=-1,\quad \xi_C'''=0,\quad
\nu'=-e/s_{a,t}^2,\quad \nu''=-1/s_{a,t}^2,\quad \nu'''=0,\quad
f'_{a,t}=a+2tz,\quad f''_{a,t}=2t,\quad f'''_{a,t}=0.
\]
After substitution, the right-hand side of the excluded identity is
\[
-\frac{2t}{f'_{a,t}(z)}+\frac{4t f'_{a,t}(z)}{s_{a,t}^2}
+\frac{4t^2\{y-f_{a,t}(z)\}}{s_{a,t}^2f'_{a,t}(z)}.
\]
For \(t>0\) and any \(z\) with \(f'_{a,t}(z)\ne0\), its coefficient of \(y\) is nonzero, whereas the left-hand side is independent of \(y\).  The identity therefore cannot hold for all admissible \((z,y)\), which verifies Assumption \(\mathrm{ass:restricted\text{-}anm}\).  All eight member properties in Definition \(\mathrm{def:base\text{-}model}\) have now been checked.  Coordinate swapping preserves these properties, and Assumption \(\mathrm{ass:pair\text{-}product}\) together with Definition \(\mathrm{def:hermite\text{-}chart}\) shows that every \(P_{a,t,\theta}\) in the asserted positive chart belongs to the same \(\mathcal M_d(3,L,c_0,C_0)\).

We calculate the local residual margin.  Put \(\tau^2=1-a^2\), let \(p_{a,t}\) be the density of \((Z,Y_t)\), and differentiate at \(t=0\).  Since \(\left.\partial_t s_{a,t}\right|_{t=0}=0\), the density score is
\[
S_a(z,y):=\left.\partial_t\log p_{a,t}(z,y)\right|_{t=0}
=\frac{(z^2-1)(y-az)}{\tau^2}.
\]
At \(t=0\), conditional Gaussian regression gives
\[
Z\mid Y_0=y=ay+U,\qquad U\sim N(0,\tau^2).
\]
Let \(m_t(y)=\mathbb E[Z\mid Y_t=y]\) and \(v_t(y)=\operatorname{Var}(Z\mid Y_t=y)\).  Differentiating their numerator and denominator integrals and using the preceding score gives
\[
\begin{aligned}
\dot m_0(y)
&=\mathbb E[(Z-ay)S_a(Z,y)\mid Y_0=y]
=a(2-3a^2)(y^2-1),\\
\dot v_0(y)
&=\mathbb E[\{(Z-ay)^2-\tau^2\}S_a(Z,y)\mid Y_0=y]
=2\tau^2(1-3a^2)y.
\end{aligned}
\]
These equalities follow by inserting \(Z=ay+U\) and the Gaussian moments \(\mathbb E U^2=\tau^2\) and \(\mathbb E U^4=3\tau^4\).  Every differentiated integrand is a polynomial times a Gaussian density whose covariance eigenvalues are bounded away from zero uniformly on \(A\); hence the differentiations and their remainders are uniform on \(A\).

Define the reverse residual \(R_t=Z-m_t(Y_t)\) and the signed measure
\[
\Delta_{a,t}=P_{Y_t,R_t}-P_{Y_t}\otimes P_{R_t}.
\]
At \(t=0\), \(R_0=Z-aY_0\) is independent of \(Y_0\), so \(\Delta_{a,0}=0\).  Since \(\mathbb E Y_t=0\) for every \(t\) and \(\mathbb E[R_t\mid Y_t]=0\),
\[
\int yr^2\,d\Delta_{a,t}(y,r)=\operatorname{Cov}(Y_t,R_t^2)=\mathbb E\{Y_tv_t(Y_t)\}.
\]
The derivative of the term involving the constant \(v_0=\tau^2\) is \(\tau^2\partial_t\mathbb E Y_t=0\).  Therefore
\[
\left.\partial_t\int yr^2\,d\Delta_{a,t}\right|_{t=0}
=\mathbb E\{Y_0\dot v_0(Y_0)\}
=2\tau^2(1-3a^2)\ne0,\qquad a\in A.
\]
Gaussian domination justifies the same identity after truncating \(yr^2\) and passing to the limit.  Consequently the signed-measure derivative \(\dot\Delta_{a,0}\) is nonzero.

The product Gaussian kernel is injective on finite signed measures.  Indeed, the elementary Fourier representation of a Gaussian gives, for every finite signed measure \(\nu\) on \(\mathbb R^2\),
\[
\left\|\int (k\otimes\ell)(z,\cdot)\,d\nu(z)\right\|_{\mathcal H_{k\otimes\ell}}^2
=c_{\sigma_K}\int_{\mathbb R^2}|\widehat\nu(\omega)|^2
e^{-\sigma_K^2\|\omega\|^2/2}\,d\omega,\qquad c_{\sigma_K}>0.
\]
If the right-hand side were zero, continuity would give \(\widehat\nu=0\) everywhere, and Fourier inversion against smooth compactly supported functions would give \(\nu=0\).  Applying this to \(\dot\Delta_{a,0}\) proves
\[
\kappa_H(a):=\left\|\mu_{\dot\Delta_{a,0}}\right\|_{\mathcal H_{k\otimes\ell}}^2>0,\qquad a\in A.
\]
Twice differentiating the bounded Gaussian-kernel features under the common Gaussian envelope yields, uniformly for \(a\in A\),
\[
\mu_{\Delta_{a,t}}=t\mu_{\dot\Delta_{a,0}}+O_{\mathcal H}(t^2)
\]
and hence
\[
H_-(a,t)=\|\mu_{\Delta_{a,t}}\|_{\mathcal H_{k\otimes\ell}}^2
=\kappa_H(a)t^2+O(t^3).
\]
The same dominated calculation makes \(a\mapsto\mu_{\dot\Delta_{a,0}}\) continuous.  Compactness of \(A\) therefore gives
\[
0<\kappa_-:=\inf_{a\in A}\kappa_H(a)\le
\sup_{a\in A}\kappa_H(a)=:\kappa_+<\infty.
\]

We next calculate the correcting gain from Definition \(\mathrm{def:gene\text{-}gain}\).  In the true direction,
\[
m_+(Z)=aZ+t(Z^2-1),\qquad q_+(a,t)=a^2+2t^2.
\]
At \(t=0\), \(m_0(Y_0)=aY_0\), so \(q_-(a,0)=a^2\).  The function \(\dot m_0(y)\) displayed above is even, \(a y\) is odd, and the marginal score \(\mathbb E[S_a(Z,Y_0)\mid Y_0=y]\) is odd.  Differentiating \(\mathbb E\{m_t(Y_t)^2\}\) thus gives zero at \(t=0\).  Uniform second differentiability gives
\[
q_-(a,t)=a^2+O(t^2).
\]
For \(t>0\), Lemma \(\mathrm{lem:population\text{-}direction\text{-}margin}\) gives \(H_+(a,t)=0\) and \(H_-(a,t)>0\).  The oracle-indicator branches in Definition \(\mathrm{def:gene\text{-}gain}\) consequently give
\[
g(a,t)=q_+(a,t)-\frac12q_-(a,t)=\frac{a^2}{2}+O(t^2),
\]
uniformly on \(A\).  A second differentiation of the same conditional Gaussian integrals also gives the more precise audit formula
\[
q_-(a,t)=a^2+2a^2(2-3a^2)^2t^2+O(t^3),\qquad
g(a,t)=\frac{a^2}{2}+c_g(a)t^2+O(t^3),
\]
where
\[
c_g(a)=2-a^2(2-3a^2)^2>0\qquad(0<a<1/2).
\]

Put \(w=t^2\) and extend the positive-branch map to its boundary by
\[
\Psi(a,w)=\bigl(g(a,\sqrt w),H_-(a,\sqrt w)\bigr),\qquad
\Psi(a,0)=(a^2/2,0).
\]
The preceding dominated expansions show that \(\Psi\) is one-sided continuously differentiable near \(A\times\{0\}\), with
\[
D\Psi(a,0)=\begin{pmatrix}
a&\partial_w g(a,0)\\
0&\kappa_H(a)
\end{pmatrix},\qquad
\det D\Psi(a,0)=a\kappa_H(a)\ge a_-\kappa_->0.
\]
Continuity permits a further decrease of \(t_0\) so that the determinant stays bounded away from zero throughout \(A\times[0,t_0^2]\).  Let \(a_0=(a_-+a_+)/2\) and \(g_0=a_0^2/2\).  To make the boundary argument explicit, reflect the map across \(w=0\) by setting
\[
\widetilde\Psi(a,w)=
\begin{cases}
\Psi(a,w),&w\ge0,\\
2\Psi(a,0)-\Psi(a,-w),&w<0.
\end{cases}
\]
The one-sided continuous derivative proved above makes \(\widetilde\Psi\) continuously differentiable near \((a_0,0)\), with the same nonsingular derivative there.  Put \(A_0=D\Psi(a_0,0)\), and choose a closed ball \(B\) centered at \(z_0=(a_0,0)\) on which
\[
\left\|A_0^{-1}\{D\widetilde\Psi(z)-A_0\}\right\|_{\mathrm{op}}<1/2.
\]
For an image point \(y\), define
\[
T_y(z)=z+A_0^{-1}\{y-\widetilde\Psi(z)\}.
\]
The mean-value identity and the preceding derivative bound show that \(T_y\) is \(1/2\)-Lipschitz on \(B\).  If \(y\) is close enough to \((g_0,0)\), then \(\|T_y(z_0)-z_0\|\) is at most one half of the radius of \(B\), so \(T_y(B)\subseteq B\).  Starting from \(z_0\), its iterates satisfy
\[
\|T_y^{k+1}(z_0)-T_y^k(z_0)\|\le2^{-k}\|T_y(z_0)-z_0\|.
\]
They are therefore Cauchy; their limit \(z(y)\in B\) satisfies \(T_y(z(y))=z(y)\), equivalently \(\widetilde\Psi(z(y))=y\).  The same Lipschitz bound proves uniqueness in \(B\).  At \(h=0\), this unique inverse is \((\sqrt{2g},0)\).  Along this boundary, implicit differentiation gives
\[
\left.\partial_h w(g,h)\right|_{h=0}=\kappa_H(\sqrt{2g})^{-1}>0.
\]
Compactness after shrinking the \(g\)-interval makes this derivative uniformly positive.  Hence \(w(g,h)>0\) for all sufficiently small \(h>0\), so the inverse then belongs to the original half-chart rather than only to its reflection.

Choose \(0<\underline g<g_0<\overline g<1\), \(h_0>0\), and \(A_H>1\) so that
\[
[\underline g,\overline g]\times[0,h_0]\subseteq\Psi(A\times[0,t_0^2]),
\]
and set \(\lambda_0=h_0/A_H\).  Then, for every \(0<\lambda\le\lambda_0\),
\[
[\underline g,\overline g]\times[\lambda,A_H\lambda]\subseteq
\Psi(A\times(0,t_0^2]).
\]

Finally choose the one-pair preimages independently for all \(j\in I\).  Product factorization makes the \(j\)-th gain and two HSIC coordinates depend only on \((a_j,t_j)\).  The true-direction HSIC is exactly zero, and Definition \(\mathrm{def:orientation\text{-}permutation}\) places that zero and the wrong-direction margin in the advertised blocks.  Thus
\[
[\underline g,\overline g]^d\times\{0\}^d\times[\lambda,A_H\lambda]^d
\subseteq
\left\{\Pi_\theta\bigl(\Gamma_d(P_{a,t,\theta},\theta)\bigr):
a\in[a_-,a_+]^d,\ 0<t_j\le t_0\right\}.
\]
Every law on the right was proved above to belong to the base model, so Definition \(\mathrm{def:feasible\text{-}set}\) gives the final inclusion in \(\Pi_\theta(\mathfrak F_d)\).  This proves all assertions, with the stronger selection \(a_+<1/2\).
\end{proof}
\par\smallskip\noindent \textit{Justification.} This unconditional feasible-image theorem and its orientation-indexed coordinate placement are the field headline: the minimax margins are populated by legal positive-density ANMs rather than imposed coordinatewise. \textit{Gap.} Li et al. (2025, Equation (6) and Theorem 4.3) assume favorable score ordering, and Peters et al. (2014) prove only qualitative identification; neither gives a local box theorem for their coupled population functionals. \textit{Consumer.} Uniform lower and upper bounds can be stated over one nonempty causal class, and benchmark designers can tune signal strength without leaving that class.

\begin{proposition}[prop:mirror-information]\label{prop:mirror-information}
Let \(a_0=(a_-+a_+)/2\).  After decreasing \(\lambda_0\) if necessary, for every \(0<\lambda\le\lambda_0\) there is \(t=t(\lambda)\in(0,t_0]\) such that \(t^2\asymp\lambda\), the two coordinate-swapped orientation laws \(P^+_{a_0,t}\) and \(P^-_{a_0,t}\) both lie in \(\mathcal M_1(\lambda)\), and constants \(0<c<C<\infty\), independent of \(\lambda\), satisfy
\[
c\lambda\le h^2(P^+_{a_0,t},P^-_{a_0,t})\le
\operatorname{KL}(P^+_{a_0,t}\Vert P^-_{a_0,t})\le C\lambda.
\]
\end{proposition}
\begin{proof}
Use the constants selected in Theorem \(\mathrm{thm:attainable\text{-}margin\text{-}rectangle}\), and put \(a_0=(a_-+a_+)/2\).  Its proof gives, at this fixed interior point,
\[
H_-(a_0,t)=\kappa_H(a_0)t^2+O(t^3),\qquad \kappa_H(a_0)>0,\qquad
g(a_0,t)=a_0^2/2+O(t^2).
\]
The map \(G(w)=H_-(a_0,\sqrt w)\) is one-sided continuously differentiable at zero with \(G(0)=0\) and \(G'(0)=\kappa_H(a_0)>0\).  Choose \(w_0>0\) so that
\[
\kappa_H(a_0)/2\le G'(w)\le3\kappa_H(a_0)/2,\qquad 0\le w\le w_0.
\]
Thus \(G\) is strictly increasing and continuous on this interval.  For every \(0<\lambda\le G(w_0)\), the intermediate-value property gives a unique \(w(\lambda)\in(0,w_0]\) with \(G(w(\lambda))=\lambda\), and integration of the derivative bounds gives
\[
\frac{2\lambda}{3\kappa_H(a_0)}\le w(\lambda)\le
\frac{2\lambda}{\kappa_H(a_0)}.
\]
Taking \(t(\lambda)=\sqrt{w(\lambda)}\) therefore supplies, after decreasing \(\lambda_0\), a continuous choice \(t=t(\lambda)\in(0,t_0]\) such that
\[
H_-(a_0,t(\lambda))=\lambda,\qquad
c_t\lambda\le t(\lambda)^2\le C_t\lambda
\]
for fixed \(0<c_t<C_t<\infty\).  The navigation interval in the theorem was chosen with \(a_0^2/2\) in its interior, so a further decrease of \(\lambda_0\) makes
\[
\underline g\le g(a_0,t(\lambda))\le\overline g.
\]
The true-direction HSIC is zero by Lemma \(\mathrm{lem:population\text{-}direction\text{-}margin}\).  Coordinate swapping interchanges the orientation labels but preserves the gain and the wrong-direction HSIC value.  Hence both \(P^+_{a_0,t(\lambda)}\) and \(P^-_{a_0,t(\lambda)}\) satisfy Definition \(\mathrm{def:separated\text{-}model}\).

It remains to compare information with \(t^2\).  With \(\tau^2=1-a_0^2\), both oriented density paths equal the centered bivariate Gaussian law \(P_{a_0,0}\) at \(t=0\).  Their respective density scores are
\[
S^+_{a_0}(x,y)=\frac{(x^2-1)(y-a_0x)}{\tau^2},\qquad
S^-_{a_0}(x,y)=\frac{(y^2-1)(x-a_0y)}{\tau^2}.
\]
Thus their score difference is
\[
D_{a_0}(x,y)=\tau^{-2}\{(x^2-1)(y-a_0x)-(y^2-1)(x-a_0y)\}.
\]
This polynomial is not zero: at \(y=0\), it equals
\[
D_{a_0}(x,0)=\tau^{-2}x\{1+a_0-a_0x^2\}.
\]
The nonsingular Gaussian base law consequently satisfies
\[
0<J_0:=\mathbb E_{a_0,0}D_{a_0}(X,Y)^2<\infty.
\]

The conditional-Gaussian location-scale paths are differentiable in quadratic mean at \(t=0\): after writing each density relative to \(P_{a_0,0}\), the explicit Gaussian likelihood ratios and their polynomial scores give an \(L^2(P_{a_0,0})\) Taylor remainder uniform for small \(|t|\).  Hence
\[
\sqrt{p^+_{a_0,t}}-\sqrt{p^-_{a_0,t}}
=\frac t2\sqrt{p_{a_0,0}}D_{a_0}+O_{L^2}(t^2).
\]
Consequently, under the convention \(h^2(P,Q)=\int(\sqrt p-\sqrt q)^2\),
\[
h^2(P^+_{a_0,t},P^-_{a_0,t})=\frac{J_0}{4}t^2+O(t^3).
\]
The same dominated likelihood-ratio expansion, using that both scores have mean zero, yields
\[
\operatorname{KL}(P^+_{a_0,t}\Vert P^-_{a_0,t})=\frac{J_0}{2}t^2+O(t^3).
\]
Finally, if \(\rho(P,Q)=\int\sqrt{dP\,dQ}\), Jensen's inequality gives
\[
\operatorname{KL}(P\Vert Q)\ge-2\log\rho(P,Q)\ge2\{1-\rho(P,Q)\}=h^2(P,Q).
\]
Combining these displays with \(t(\lambda)^2\asymp\lambda\), and decreasing \(\lambda_0\) once more to absorb the remainders, proves the stated two-sided Hellinger and upper Kullback--Leibler bounds.
\end{proof}
\par\smallskip\noindent \textit{Justification.} The mirror calculation links the causal margin \(\lambda\) to one-sample testing information without introducing an unrelated divergence assumption. \textit{Gap.} Acharya et al. (2023, Theorems 1.1--1.5) derive intervention--observation tradeoffs for discrete causal-structure identification, and Jin (2026, arXiv:2608.15840v1, Theorem 3.1) derives a bivariate linear-LiNGAM frontier from edge, non-Gaussianity, and scale margins; neither constructs opposite-orientation nonlinear restricted-ANM laws whose fixed-kernel residual-HSIC margin and testing information are both of order \(\lambda\). \textit{Consumer.} The proposition is the legal one-coordinate experiment needed for the product-testing converse.

\begin{theorem}[thm:product-testing-converse]\label{thm:product-testing-converse}
There are universal constants \(c_1,c_2>0\) such that, for all \(d\ge1\), \(0<\alpha\le1/4\), and \(0<\lambda\le\lambda_0\), every measurable estimator \(\widetilde\theta\) satisfies \[n\le c_1\frac{\log((d+1)/\alpha)}{\lambda}\quad\Longrightarrow\quad\sup_{(P,\theta)\in\mathcal M_d(\lambda)}P^{\otimes n}(\widetilde\theta\ne\theta)>\alpha.\] Moreover, any random set \(\widetilde{\mathcal C}\subseteq\{-1,+1\}^I\) with uniform coverage at least \(1-\alpha\) obeys \[\sup_{(P,\theta)\in\mathcal M_d(\lambda)}P^{\otimes n}(|\widetilde{\mathcal C}|>1)\ge c_2\] in the same subcritical regime after a universal adjustment of \(c_1\).
\end{theorem}
\begin{proof}
Let \(Q_+=P^+_{a_0,t(\lambda)}\) and \(Q_-=P^-_{a_0,t(\lambda)}\) be the calibrated mirror laws from Proposition \(\mathrm{prop:mirror\text{-}information}\).  For each \(\vartheta\in\{-1,+1\}^d\), define
\[
Q_\vartheta=\bigotimes_{j=1}^d Q_{\vartheta_j}.
\]
The proposition gives the one-pair base-model and margin properties, while Assumption \(\mathrm{ass:pair\text{-}product}\) gives the product property.  Thus every \((Q_\vartheta,\vartheta)\) is a member of the same \(\mathcal M_d(\lambda)\).  It is therefore enough to lower-bound risk on this finite product subexperiment.

We first derive the binary testing inequality used below.  For laws \(R,S\), let
\[
\rho(R,S)=\int\sqrt{dR\,dS}.
\]
Jensen's inequality and Cauchy--Schwarz respectively give
\[
\rho(R,S)\ge \exp\{-\operatorname{KL}(R\Vert S)/2\},\qquad
\|R-S\|_{\mathrm{TV}}\le\sqrt{1-\rho(R,S)^2}.
\]
Hence the equal-prior Bayes error \(e(R,S)\) satisfies
\[
e(R,S)=\frac{1-\|R-S\|_{\mathrm{TV}}}{2}
\ge\frac{1-\sqrt{1-\exp\{-\operatorname{KL}(R\Vert S)\}}}{2}
\ge\frac14\exp\{-\operatorname{KL}(R\Vert S)\}.
\]
By the mirror proposition, for the \(n\)-sample one-coordinate problem there is a fixed numerical \(K<\infty\) such that
\[
e_n:=e(Q_+^{\otimes n},Q_-^{\otimes n})\ge
\frac{1-\sqrt{1-e^{-Kn\lambda}}}{2}\ge\frac14e^{-Kn\lambda}.
\]

Under the uniform prior on \(\vartheta\), the prior and likelihood factor over the \(d\) coordinates.  The posterior therefore factors, and the Bayes-optimal exact-vector rule is the product of the coordinatewise binary Bayes rules.  Every estimator consequently obeys
\[
\sup_{\vartheta}Q_\vartheta^{\otimes n}(\widetilde\theta\ne\vartheta)
\ge1-(1-e_n)^d.
\]
Put \(R_\alpha=(d+1)/\alpha\).  Choose \(c_1>0\) so that \(\eta:=Kc_1\le1/64\).  If \(n\le c_1\lambda^{-1}\log R_\alpha\), then \(e^{-Kn\lambda}\ge R_\alpha^{-\eta}\).  If \(R_\alpha\le128\), then
\[
R_\alpha^{-\eta}\ge128^{-1/64}>3/4,\qquad e_n>1/4\ge\alpha,
\]
so the vector error is strictly larger than \(\alpha\).  If \(R_\alpha>128\), put \(u=d/\alpha\).  Since \(d+1\le2d\), one has \(R_\alpha\le2u\) and \(u>64\).  Therefore
\[
\frac{de_n}{\alpha}\ge\frac{u}{4R_\alpha^\eta}\ge
\frac{u^{1-\eta}}{4\,2^\eta}>2.
\]
If \(de_n\le1\), then
\[
1-(1-e_n)^d\ge1-e^{-de_n}\ge de_n/2>\alpha;
\]
if \(de_n>1\), the same quantity exceeds \(1-e^{-1}>1/4\ge\alpha\).  This proves the exact-orientation assertion after the stated universal choice of \(c_1\).

We now prove the confidence-set assertion.  Fix \(c_2=1/32\), and decrease \(c_1\) below when necessary.  First suppose
\[
\log(d+1)\ge\log(1/\alpha).
\]
Then \(d\ge3\), \(\log R_\alpha\le2\log(d+1)\), and, after requiring \(2Kc_1\le1/64\),
\[
e_n\ge\frac1{4(d+1)^{1/64}}.
\]
It follows that
\[
1-(1-e_n)^d\ge1-e^{-de_n}\ge
1-\exp\left\{-\frac{3}{4\cdot4^{1/64}}\right\}>\frac12>\alpha+c_2.
\]
If a random set \(\widetilde{\mathcal C}\) had coverage at least \(1-\alpha\) and \(Q_\vartheta^{\otimes n}(|\widetilde{\mathcal C}|>1)<c_2\) for every \(\vartheta\), output its unique element when it is a singleton and an arbitrary fixed vector otherwise.  For every \(\vartheta\), the resulting estimator would satisfy
\[
Q_\vartheta^{\otimes n}(\widetilde\theta\ne\vartheta)\le
Q_\vartheta^{\otimes n}\{\vartheta\notin\widetilde{\mathcal C}\}
+Q_\vartheta^{\otimes n}\{|\widetilde{\mathcal C}|>1\}<\alpha+c_2,
\]
contradicting the preceding Bayes lower bound.

It remains to consider
\[
\log(1/\alpha)>\log(d+1).
\]
Take two product laws \(R_+\) and \(R_-\) whose orientation vectors \(\vartheta^+\) and \(\vartheta^-\) differ only in the first coordinate.  Suppose again, for contradiction, that the ambiguity probability is less than \(c_2\) under every law, and define
\[
B=\{\widetilde{\mathcal C}=\{\vartheta^+\}\}.
\]
Coverage and the ambiguity bound imply
\[
R_+^{\otimes n}(B)\ge1-\alpha-c_2,\qquad
R_-^{\otimes n}(B)\le\alpha.
\]
Binary data processing for Kullback--Leibler divergence, obtained directly by applying the log-sum inequality to \(B\) and \(B^c\), gives
\[
\begin{aligned}
n\operatorname{KL}(R_+\Vert R_-)
&\ge \operatorname{kl}(1-\alpha-c_2,\alpha)\\
&\ge(1-\alpha-c_2)\log(1/\alpha)-\log2\\
&\ge\frac7{32}\log(1/\alpha).
\end{aligned}
\]
The last inequality uses \(1-\alpha-c_2\ge23/32\), \(\alpha\le1/4\), and hence \(\log2\le\frac12\log(1/\alpha)\).  On the other hand, the two laws differ in one factor only, so the mirror information bound and the subcritical sample-size condition give
\[
n\operatorname{KL}(R_+\Vert R_-)\le Kn\lambda
\le Kc_1\log R_\alpha<2Kc_1\log(1/\alpha).
\]
Choosing \(2Kc_1<7/32\) is a contradiction.  Thus in both regimes some member of \(\mathcal M_d(\lambda)\) has ambiguity probability at least \(c_2\).  All constants are numerical once the fixed chart and regularity envelope of Theorem \(\mathrm{thm:attainable\text{-}margin\text{-}rectangle}\) are fixed, completing the proof.
\end{proof}
\par\smallskip\noindent \textit{Justification.} This supplies the all-procedure necessity of both the confidence and dimension logarithms under exact simultaneous loss. \textit{Gap.} Existing minimax HSIC results treat one observed-variable independence problem. Acharya et al. (2023) use interventions and a discrete total-variation causal-structure margin, while Jin (2026, Theorem 3.1) uses bivariate linear LiNGAM primitive margins; none gives the \(\log d\) exact-recovery lower bound through legal nonlinear residual-HSIC product mirrors. \textit{Consumer.} No causal-discovery algorithm can claim a uniformly earlier singleton transition on the stated supplied-pair experiment.

\begin{theorem}[thm:uniform-minimax-frontier]\label{thm:uniform-minimax-frontier}
Fix the sub-half Hermite chart selected in Lemma \(\mathrm{lem:subhalf\text{-}chart\text{-}calibration}\).  There are constants \(0<c<C<\infty\), depending only on this fixed chart and on the fixed Gaussian kernels, such that for every \(d\ge1\), \(0<\alpha\le1/4\), and \(0<\lambda\le\lambda_H\),
\[
n\lambda\ge C\log\frac{d+1}{\alpha}
\quad\Longrightarrow\quad
\sup_{(P,\theta)\in\mathcal M_d^{\mathrm H}(\lambda)}
P^{\otimes n}(\widehat\theta_n^{\mathrm H}\ne\theta)\le\alpha/2,
\]
where \(\widehat\theta_n^{\mathrm H}\) is the explicit median-of-means rule in Definition \(\mathrm{def:chart\text{-}moment\text{-}rule}\).  Every measurable estimator has worst-case error greater than \(\alpha\) when \(n\lambda\le c\log((d+1)/\alpha)\).  Consequently,
\[
c\frac{\log((d+1)/\alpha)}{\lambda}
\le N_{\mathrm H}^*(d,\alpha,\lambda)
\le C\frac{\log((d+1)/\alpha)}{\lambda}.
\]
Since \(\mathcal M_d^{\mathrm H}(\lambda)\subset\mathcal M_d(\lambda)\), the same lower bound transfers to the full-class frontier \(N^*(d,\alpha,\lambda)\), but no full-class upper bound is asserted.
\end{theorem}
\begin{proof}
Put
\[
L_{d,\alpha}=\log\frac{d+1}{\alpha},\qquad
K=\left\lceil 8\log\frac{2d}{\alpha}\right\rceil,
\qquad q=\left\lfloor\frac nK\right\rfloor,
\qquad \rho=\sqrt{\frac{8V_\psi}{q}},
\]
as in Definition \(\mathrm{def:chart\text{-}moment\text{-}rule}\).  Fix
\((P,\theta)\in\mathcal M_d^{\mathrm H}(\lambda)\), and define
\[
\mu_j=\mathbb E_P\psi(X_j,Y_j).
\]
Lemma \(\mathrm{lem:subhalf\text{-}chart\text{-}calibration}\) gives, for every \(j\in I\),
\[
\mu_j=\theta_j\,2t_j(1-2a_j),
\qquad
\operatorname{Var}_P\{\psi(X_j,Y_j)\}\le V_\psi,
\qquad
H_{j,-\theta_j}\le C_Ht_j^2.
\tag{1}
\]
In particular, Definition \(\mathrm{def:subhalf\text{-}hermite\text{-}class}\) imposes \(t_j>0\) and \(a_j\le a_+<1/2\), so
\[
\operatorname{sign}(\mu_j)=\theta_j.
\tag{2}
\]

Suppose first that \(n\ge2K\).  For a block \(A_k\), let
\[
\overline\psi_{jk}=q^{-1}\sum_{i\in A_k}\psi_{ij}.
\]
The independent-sample model and (1) imply that this block mean has mean \(\mu_j\) and variance at most \(V_\psi/q\).  Chebyshev's inequality, obtained directly by applying Markov's inequality to \((\overline\psi_{jk}-\mu_j)^2\), gives
\[
P^{\otimes n}\bigl(|\overline\psi_{jk}-\mu_j|>\rho\bigr)
\le \frac{V_\psi}{q\rho^2}=\frac18.
\tag{3}
\]
The \(K\) block means are independent because the observation-index blocks are disjoint.  If their lower median differs from \(\mu_j\) by more than \(\rho\), at least \(K/2\) block means satisfy the event in (3).  A union bound over subsets of blocks therefore yields
\[
\begin{aligned}
P^{\otimes n}\bigl(|\widehat\mu_j-\mu_j|>\rho\bigr)
&\le \sum_{s=\lceil K/2\rceil}^{K}\binom Ks(1/8)^s\\
&\le 2^K(1/8)^{K/2}
=2^{-K/2}\\
&\le \exp\left\{-4(\log2)\log\frac{2d}{\alpha}\right\}
\le \frac{\alpha}{2d}.
\end{aligned}
\tag{4}
\]
For the last inequality, \(K\ge8\log(2d/\alpha)\), \(4\log2>1\), and \(\alpha/(2d)<1\).  Thus a union bound over the deterministic set \(I\) gives
\[
P^{\otimes n}(\mathcal E)\ge1-\alpha/2,
\qquad
\mathcal E=\left\{\max_{j\in I}|\widehat\mu_j-\mu_j|\le\rho\right\}.
\tag{5}
\]
This probability bound is uniform over the chart class because the variance bound in (1) is uniform.

Membership in \(\mathcal M_d^{\mathrm H}(\lambda)\) entails membership in \(\mathcal M_d(\lambda)\), hence \(H_{j,-\theta_j}\ge\lambda\).  Combining this fact with (1) gives
\[
t_j^2\ge\frac{\lambda}{C_H},
\qquad
|\mu_j|^2=4t_j^2(1-2a_j)^2
\ge \frac{4(1-2a_+)^2}{C_H}\lambda.
\tag{6}
\]
When \(n\ge2K\), \(q=\lfloor n/K\rfloor\ge n/(2K)\), and therefore
\[
\rho^2\le\frac{16V_\psi K}{n}.
\tag{7}
\]
For \(d\ge1\) and \(0<\alpha\le1/4\),
\[
\log\frac{2d}{\alpha}\le2L_{d,\alpha},
\qquad
K\le9\log\frac{2d}{\alpha}\le18L_{d,\alpha}.
\tag{8}
\]
Choose \(C\), depending only on \(V_\psi\), \(C_H\), and the fixed endpoint \(a_+\), so that
\[
C>\max\left\{36,\frac{288V_\psi C_H}{(1-2a_+)^2}\right\}.
\tag{9}
\]
Because \(0<\lambda\le\lambda_H\le1\), the inequality \(n\lambda\ge CL_{d,\alpha}\), together with (8)--(9), implies \(n\ge2K\).  Equations (6)--(9) also imply
\[
4\rho^2<\min_{j\in I}|\mu_j|^2.
\tag{10}
\]
Indeed, (7)--(8) give \(4\rho^2\le1152V_\psi L_{d,\alpha}/n\le1152V_\psi\lambda/C\), while (6) bounds the right-hand side of (10) below by \(4(1-2a_+)^2\lambda/C_H\), and (9) compares these constants.

On \(\mathcal E\), (10) implies that each interval \([\widehat\mu_j-\rho,\widehat\mu_j+\rho]\) lies strictly in the open half-line with sign \(\theta_j\).  Definition \(\mathrm{def:chart\text{-}moment\text{-}rule}\) then gives \(B_j=\{\theta_j\}\) for every \(j\), and consequently
\[
\widehat{\mathcal C}_n^{\mathrm H}=\{\theta\},
\qquad
\widehat\theta_n^{\mathrm H}=\theta
\quad\text{on }\mathcal E.
\]
Together with (5), this proves
\[
\sup_{(P,\theta)\in\mathcal M_d^{\mathrm H}(\lambda)}
P^{\otimes n}(\widehat\theta_n^{\mathrm H}\ne\theta)
\le\alpha/2.
\tag{11}
\]

Lemma \(\mathrm{lem:chart\text{-}product\text{-}testing\text{-}lower}\) supplies \(c_L>0\) such that every measurable estimator has worst-case error greater than \(\alpha\) whenever \(n\lambda\le c_LL_{d,\alpha}\).  Combining that lower bound with (11) and Definition \(\mathrm{def:chart\text{-}minimax\text{-}frontier}\), and absorbing the integer ceiling into the constants using \(L_{d,\alpha}\ge\log8\) and \(\lambda\le1\), yields constants \(0<c<C<\infty\) for which
\[
c\frac{L_{d,\alpha}}{\lambda}
\le N_{\mathrm H}^*(d,\alpha,\lambda)
\le C\frac{L_{d,\alpha}}{\lambda}.
\]
Finally, Definition \(\mathrm{def:subhalf\text{-}hermite\text{-}class}\) gives \(\mathcal M_d^{\mathrm H}(\lambda)\subseteq\mathcal M_d(\lambda)\).  The supremum risk over \(\mathcal M_d(\lambda)\) is therefore no smaller than the supremum over its chart subclass, so the same lower bound holds for \(N^*(d,\alpha,\lambda)\); this also agrees with Theorem \(\mathrm{thm:product\text{-}testing\text{-}converse}\).  Lemma \(\mathrm{lem:peters\text{-}disjoint\text{-}pair\text{-}specialization}\), applied to the member conditions in Definition \(\mathrm{def:base\text{-}model}\), ensures that \(\theta\) is the observationally identified causal orientation on the disjoint-edge restricted-ANM domain.  Neither the chart estimator nor its upper bound is transferred to the larger class.
\end{proof}
\par\smallskip\noindent \textit{Justification.} The theorem gives an unconditional all-procedure matched frontier on the explicit sub-half Hermite chart: the median-of-means observable-moment estimator attains the rate and the product-testing lemma proves that no measurable estimator improves its order. Class containment transfers only the lower bound to the larger separated restricted-ANM class. \textit{Gap.} Pointwise consistency of regression--independence methods does not provide a finite-sample, familywise, matched exact-recovery rate for supplied independent ANM pairs; no cited result supplies the chart-uniform frontier proved here. \textit{Consumer.} Practitioners working within the calibrated Hermite chart obtain the explicit sample-size rule \(n\lambda\asymp\log((d+1)/\alpha)\); outside that chart the theorem supplies a hardness benchmark only.

\begin{theorem}[thm:simultaneous-graph-confidence-set]\label{thm:simultaneous-graph-confidence-set}
For every \(d\ge1\), \(0<\alpha\le1/4\), and \(0<\lambda\le\lambda_H\), put \(K=\lceil8\log(2d/\alpha)\rceil\).  For every \(n\ge2K\), the moment confidence set \(\widehat{\mathcal C}_n^{\mathrm H}\) in Definition \(\mathrm{def:chart\text{-}moment\text{-}rule}\) is measurable, uses exactly \(d\) observable moment queries, has arithmetic cost \(O(dn)\), and satisfies
\[
\inf_{(P,\theta)\in\mathcal M_d^{\mathrm H}(\lambda)}
P^{\otimes n}\{\theta\in\widehat{\mathcal C}_n^{\mathrm H}\}
\ge1-\alpha/2.
\]
For the constant \(C\) in Theorem \(\mathrm{thm:uniform\text{-}minimax\text{-}frontier}\),
\[
n\lambda\ge C\log\frac{d+1}{\alpha}
\quad\Longrightarrow\quad
\inf_{(P,\theta)\in\mathcal M_d^{\mathrm H}(\lambda)}
P^{\otimes n}\{\widehat{\mathcal C}_n^{\mathrm H}=\{\theta\}\}
\ge1-\alpha/2.
\]
Conversely, below the constant-factor threshold in Lemma \(\mathrm{lem:chart\text{-}product\text{-}testing\text{-}lower}\), no random set over the same chart class can have both uniform coverage and uniform singleton probability at least \(1-\alpha/2\).
\end{theorem}
\begin{proof}
All block boundaries in Definition \(\mathrm{def:chart\text{-}moment\text{-}rule}\) are deterministic.  Each block average is a Borel function of the sample, a finite lower order statistic is Borel measurable, and every condition defining \(B_j\) is a Borel inequality.  Since \(\{-1,+1\}^I\) is finite, the Cartesian product \(\widehat{\mathcal C}_n^{\mathrm H}=\prod_{j\in I}B_j\) is a measurable random set.  The rule evaluates the fixed observable functional \(\psi\) once for each pair coordinate and observation, hence it uses exactly \(d\) coordinatewise moment queries.  Its \(nd\) evaluations and block sums require \(O(dn)\) arithmetic operations.  A deterministic linear-time selection algorithm computes each median in \(O(K)\) operations, and \(K\le n/2\) because \(n\ge2K\), so the total cost remains \(O(dn)\).

Fix \((P,\theta)\in\mathcal M_d^{\mathrm H}(\lambda)\) and write
\[
\mu_j=\mathbb E_P\psi(X_j,Y_j).
\]
The block argument in the proof of Theorem \(\mathrm{thm:uniform\text{-}minimax\text{-}frontier}\), using the uniform variance bound in Lemma \(\mathrm{lem:subhalf\text{-}chart\text{-}calibration}\), proves
\[
P^{\otimes n}\left\{\max_{j\in I}|\widehat\mu_j-\mu_j|
\le\rho_{n,d,\alpha}\right\}
\ge1-\alpha/2.
\tag{12}
\]
On the event in (12), the interval \([\widehat\mu_j-\rho_{n,d,\alpha},\widehat\mu_j+\rho_{n,d,\alpha}]\) contains \(\mu_j\).  The calibration lemma gives
\[
\mu_j=\theta_j\,2t_j(1-2a_j).
\]
Definition \(\mathrm{def:subhalf\text{-}hermite\text{-}class}\) gives \(t_j>0\) and \(a_j\le a_+<1/2\), so \(\operatorname{sign}(\mu_j)=\theta_j\).  If \(\theta_j=+1\), interval containment gives \(\widehat\mu_j+\rho_{n,d,\alpha}\ge\mu_j>0\), and hence \(+1\in B_j\).  If \(\theta_j=-1\), it gives \(\widehat\mu_j-\rho_{n,d,\alpha}\le\mu_j<0\), and hence \(-1\in B_j\).  Thus \(\theta_j\in B_j\) in both cases, and (12) proves
\[
\inf_{(P,\theta)\in\mathcal M_d^{\mathrm H}(\lambda)}
P^{\otimes n}\left\{\theta\in\widehat{\mathcal C}_n^{\mathrm H}\right\}
\ge1-\alpha/2.
\tag{13}
\]
Because \(t_j>0\), the chart mechanism is non-affine.  Lemma \(\mathrm{lem:population\text{-}direction\text{-}margin}\) therefore gives \(H_{j,\theta_j}=0<H_{j,-\theta_j}\).  Hence the sign used above recovers the causal orientation identified by the population residual-HSIC functional; it does not define a different target.

For the constant \(C\) in Theorem \(\mathrm{thm:uniform\text{-}minimax\text{-}frontier}\), the implication \(n\lambda\ge C\log((d+1)/\alpha)\) gives
\[
2\rho_{n,d,\alpha}<\min_{j\in I}|\mu_j|.
\tag{14}
\]
On the event in (12), (14) places every interval strictly in the half-line with sign \(\theta_j\).  Thus \(B_j=\{\theta_j\}\) for every \(j\), and
\[
\widehat{\mathcal C}_n^{\mathrm H}=\{\theta\}.
\]
Combining this identity with (12) proves the asserted uniform singleton probability.  Finally, Lemma \(\mathrm{lem:chart\text{-}product\text{-}testing\text{-}lower}\) states, on the same chart class and below its constant-factor threshold, that no random set can have both uniform coverage and uniform singleton probability at least \(1-\alpha/2\).  This proves the converse assertion.
\end{proof}
\par\smallskip\noindent \textit{Justification.} The confidence set is an unconditional, measurable product inversion of the observable chart moments. Its finite-sample coverage follows from the uniform variance envelope, and its singleton transition occurs at the same chart-minimax threshold without a nuisance regression event. \textit{Gap.} Existing asymptotic ordering sets do not provide this finite-sample \(d\)-query product confidence set with a matching impossibility of simultaneous coverage and singleton output below the chart threshold. \textit{Consumer.} Analysts on the explicit chart can report every orientation vector compatible with the simultaneous moment intervals, with \(O(dn)\) arithmetic cost and uniform coverage.

\begin{proposition}[prop:single-pair-reduction]\label{prop:single-pair-reduction}
When \(d=1\), Definition \(\mathrm{def:chart\text{-}moment\text{-}rule}\) gives \(\widehat{\mathcal C}_n^{\mathrm H}=B_1\): it contains \(+1\) exactly when \(\widehat\mu_1+\rho_{n,1,\alpha}\ge0\) and contains \(-1\) exactly when \(\widehat\mu_1-\rho_{n,1,\alpha}\le0\).  On the sub-half chart the population sign of \(\mathbb E\psi(X_1,Y_1)\) equals the unique zero-residual-HSIC direction, and
\[
N_{\mathrm H}^*(1,\alpha,\lambda)\asymp
\lambda^{-1}\log(2/\alpha),
\qquad 0<\alpha\le1/4,\quad 0<\lambda\le\lambda_H.
\]
\end{proposition}
\begin{proof}
When \(d=1\), one has \(I=\{1\}\), so the Cartesian product in Definition \(\mathrm{def:chart\text{-}moment\text{-}rule}\) is
\[
\widehat{\mathcal C}_n^{\mathrm H}=\prod_{j\in I}B_j=B_1.
\]
Substituting \(j=1\) into that definition gives
\[
+1\in B_1
\quad\Longleftrightarrow\quad
\widehat\mu_1+\rho_{n,1,\alpha}\ge0,
\qquad
-1\in B_1
\quad\Longleftrightarrow\quad
\widehat\mu_1-\rho_{n,1,\alpha}\le0,
\]
which is the asserted pairwise inversion.

For every member of \(\mathcal M_1^{\mathrm H}(\lambda)\), Lemma \(\mathrm{lem:subhalf\text{-}chart\text{-}calibration}\) gives
\[
\mathbb E_P\psi(X_1,Y_1)=\theta_1\,2t_1(1-2a_1).
\]
The sub-half restrictions impose \(t_1>0\) and \(a_1\le a_+<1/2\), so the sign of this expectation is \(\theta_1\).  Since \(t_1>0\), the chart mechanism is non-affine.  Lemma \(\mathrm{lem:population\text{-}direction\text{-}margin}\) then yields
\[
H_{1,\theta_1}=0<H_{1,-\theta_1},
\]
showing that the moment sign equals the unique residual-HSIC-identified causal direction.

Finally, specialize the two-sided conclusion of Theorem \(\mathrm{thm:uniform\text{-}minimax\text{-}frontier}\) to \(d=1\).  Since \(\log((d+1)/\alpha)=\log(2/\alpha)\), constants independent of \(\alpha\) and \(\lambda\) satisfy
\[
c\lambda^{-1}\log\frac2\alpha
\le N_{\mathrm H}^*(1,\alpha,\lambda)
\le C\lambda^{-1}\log\frac2\alpha,
\qquad
0<\alpha\le1/4,\quad 0<\lambda\le\lambda_H.
\]
This is the asserted \(\asymp\) relation.  The coverage interpretation of the displayed inversion follows by specializing Theorem \(\mathrm{thm:simultaneous\text{-}graph\text{-}confidence\text{-}set}\) to \(d=1\).
\end{proof}
\par\smallskip\noindent \textit{Justification.} The proposition verifies that the product construction reduces at \(d=1\) to inversion of one antisymmetric chart-moment interval and that its sign targets the unique zero-residual-HSIC causal direction. \textit{Gap.} Qualitative bivariate ANM identification and pointwise consistency do not yield the local finite-sample complexity \(\lambda^{-1}\log(2/\alpha)\) established by the chart specialization. \textit{Consumer.} Single-pair users obtain a direct implementation check and the calibrated one-pair sample-complexity benchmark.

\begin{theorem}[thm:fixed-local-polynomial-uniform-expansion-impossible]\label{thm:fixed-local-polynomial-uniform-expansion-impossible}
For the fixed four-fold local-polynomial class \(\mathcal A_n^{\mathrm{LP}}\), there exist a law \(P^\circ\in\mathcal M_1(\beta,L,c_0,C_0)\) and a constant \(\delta\in(0,1-\log 2)\) such that, for every finite \(C_m>0\), the unique admissible fit \(\widehat m\), and \(x=1\),
\[
\liminf_{n\to\infty}P^{\circ\otimes n}\!\left\{\mathcal G_n(1;C_m,\widehat m)^c\right\}\ge e^{-\delta}>2e^{-1}.
\]
The exact necessary probability requirement for the fixed local-polynomial/orthogonalization protocol is the following: for some constants \(c_u,C_H,C_m>0\), depending only on the fixed envelope and kernels, and for measurable statistics \((\widehat H^\perp_{j,b}(\widehat m):(j,b)\in\mathcal Q_d)\) and associated data-dependent radii produced within \(\mathfrak H_{\mathrm{orth}}(\mathcal A_n^{\mathrm{LP}})\), uniformly over \(P\in\mathcal M_d(\beta,L,c_0,C_0)\), every \(n\ge1\), every \(\widehat m\in\mathcal A_n^{\mathrm{LP}}\), and every \(x\ge1\) for which \(u=x+\log(2d)\le c_u n\), one must have
\[
P^{\otimes n}\!\left(
\mathcal G_n(x;C_m,\widehat m)^c
\ \cup\
\left\{\exists(j,b)\in\mathcal Q_d:
\left|\widehat H^\perp_{j,b}(\widehat m)-H_{j,b}\right|
>C_H\left(\sqrt{\frac{H_{j,b}u}{n}}+\frac{u}{n}+r_n^4\right)
\right\}
\right)\le2e^{-x}.
\]
At \(d=1\) and \(x=1\), the failure union contains \(\mathcal G_n(1;C_m,\widehat m)^c\) and the required upper bound is \(2e^{-1}\); hence the displayed requirement is impossible for every choice of such statistics and radii within this fixed protocol. This conclusion is not an impossibility result for arbitrary regression estimators, arbitrary independence statistics, or all residual-HSIC error procedures.
\end{theorem}
\begin{proof}
Set \(d=1\) and \(\theta=+1\). By Theorem \(\mathrm{thm:attainable\text{-}margin\text{-}rectangle}\), choose \(a\) in the interior of its certified chart interval and fix \(t\in(0,t_0]\). The nonlinear chart law
\[
P^\circ=P_{a,t,+}
\]
then belongs to \(\mathcal M_1(\beta,L,c_0,C_0)\). Under Definition \(\mathrm{def:hermite\text{-}chart}\), its true-direction query \(b=+1\) has
\[
C=Z,qquad E=aZ+t(Z^2-1)+s_{a,t}\xi,qquad
Z\mathrel{\perp\!\!\!\perp}\xi,qquad Z,\xi\sim\mathcal N(0,1).
\]
Thus \(\mathbb E(\xi\mid Z)=0\), and Definition \(\mathrm{def:population\text{-}regressions}\) gives
\[
m(z)=\mathbb E_{P^\circ}(E\mid Z=z)=az+t(z^2-1).
\]

Fix one of the four regression folds in Definition \(\mathrm{def:balanced\text{-}four\text{-}fold\text{-}allocation}\), and write \(m_n\) for its cardinality. Balance gives \(m_n/n\to1/4\). Put \(h=h_n\). Definition \(\mathrm{def:admissible\text{-}local\text{-}polynomial\text{-}class}\) gives
\[
h=\{\max(1,\lfloor n/4\rfloor)\}^{-1/(2\beta+1)},
\qquad m_nh\longrightarrow\infty,
\]
where the limit follows because \(\beta>1/2\), hence \(1-1/(2\beta+1)>0\).

Since \(\log2=\int_1^2z^{-1}\,dz<1\), choose
\[
0<\delta<1-\log2.
\]
Then \(e^{-\delta}>e^{-(1-\log2)}=2e^{-1}\). Let \(\varphi_0(z)=(2\pi)^{-1/2}e^{-z^2/2}\) be the standard-normal density and, for \(z\ge2h\), define
\[
F_n(z)=P^\circ\{Z\in[z-2h,z+2h]\}.
\]
The map \(F_n\) is continuous, \(F_n(z)\to0\) as \(z\to\infty\), and it is strictly decreasing on \((2h,\infty)\), because
\[
F_n'(z)=\varphi_0(z+2h)-\varphi_0(z-2h)<0.
\]
Also \(F_n(2h)=\int_0^{4h}\varphi_0(v)\,dv\sim4\varphi_0(0)h\), so \(m_nF_n(2h)\to\infty\). The intermediate value theorem therefore gives, for all sufficiently large \(n\), a unique \(z_n>2h\) such that, with
\[
J_n=[z_n-2h,z_n+2h],
\qquad m_nP^\circ(Z\in J_n)=\delta.
\]

The integral mean-value theorem supplies \(\zeta_n\in J_n\) for which
\[
4h\varphi_0(\zeta_n)=P^\circ(Z\in J_n)=\frac{\delta}{m_n}.
\]
Consequently,
\[
\zeta_n^2
=2\log\!\left(\frac{4m_nh}{\delta\sqrt{2\pi}}\right)
=\frac{4\beta}{2\beta+1}\log n+O(1).
\]
In particular, \(\zeta_n\to\infty\). Because \(|z_n-\zeta_n|\le2h\), the preceding display first implies \(z_n=O(\sqrt{\log n})\) and hence \(z_nh\to0\); it then yields
\[
|z_n^2-\zeta_n^2|
\le2h(z_n+\zeta_n)=o(1),
\qquad
z_n^2=\frac{4\beta}{2\beta+1}\log n+O(1).
\]
Since \(4\beta/(2\beta+1)<2\), whereas the same admissible-protocol definition fixes
\[
T_n^2=\frac{8\log(n\vee3)}{c_0}
\]
with \(0<c_0<1\), it follows that \(z_n+3h/2<T_n\) for all sufficiently large \(n\).

Define
\[
A_n=[z_n-h/2,z_n+h/2].
\]
For any \(v,w\in J_n\),
\[
\frac{\varphi_0(v)}{\varphi_0(w)}
=\exp\{(w^2-v^2)/2\},
\qquad
|v^2-w^2|\le8z_nh+16h^2.
\]
Because \(z_nh+h^2\to0\), the ratio of the maximum and minimum of \(\varphi_0\) on \(J_n\) tends to one. The lengths of \(A_n\) and \(J_n\) are \(h\) and \(4h\), respectively, and hence
\[
\frac{P^\circ(Z\in A_n)}{P^\circ(Z\in J_n)}\longrightarrow\frac14.
\]
Using \(P^\circ(Z\in J_n)=\delta/m_n\) and \(m_n\le n\), there is a constant \(c_A>0\) such that
\[
P^\circ(Z\in A_n)\ge\frac{c_A}{n}
\]
for all sufficiently large \(n\).

Let \(E_n\) be the event that the selected regression fold contains no predictor value in \(J_n\). The sample observations are independent under \(P^{\circ\otimes n}\), so
\[
P^{\circ\otimes n}(E_n)
=\left(1-\frac{\delta}{m_n}\right)^{m_n}.
\]
Since
\[
m_n\log\left(1-\frac{\delta}{m_n}\right)\longrightarrow-\delta,
\]
we obtain
\[
P^{\circ\otimes n}(E_n)\longrightarrow e^{-\delta}.
\]

We now evaluate the unique fit specified by Definition \(\mathrm{def:admissible\text{-}local\text{-}polynomial\text{-}class}\). If \(z\in A_n\), then \(|z|\le T_n-h\) for all sufficiently large \(n\), so the interior kernel \(K_{z,n}=K_{\mathrm{LP}}\) is used. Its support implies that a nonzero observation weight requires
\[
|Z_i-z|\le h,
\]
which implies
\[
Z_i\in[z_n-3h/2,z_n+3h/2]\subset J_n.
\]
On \(E_n\), therefore, every data-fit term in the local-polynomial objective is zero for every \(z\in A_n\). The remaining objective is
\[
n^{-2}\sum_{q=0}^{p_{\mathrm{LP}}}a_q^2,
\]
whose unique minimizer is \(a_q=0\) for every \(q\). Clipping leaves zero unchanged, and the selected foldwise true-direction fit consequently satisfies
\[
\widehat m(z)=0\qquad\text{for every }z\in A_n
\]
on \(E_n\).

Because \(t>0\), \(a>0\), and \(z_n\to\infty\), the formula \(m(z)=az+t(z^2-1)\) yields constants \(c_m,c_z>0\) such that, for all sufficiently large \(n\),
\[
\inf_{z\in A_n}|m(z)|\ge c_mz_n^2\ge c_z\log n.
\]
Combining this inequality with the probability lower bound for \(A_n\), on \(E_n\) we have
\[
\begin{aligned}
\|\widehat m-m\|_{L^4(P^\circ_C)}
&\ge\left\{\int_{A_n}|m(z)|^4\,dP^\circ_C(z)\right\}^{1/4}\\
&\ge c_E n^{-1/4}\log n
\end{aligned}
\]
for a fixed constant \(c_E>0\).

By Definition \(\mathrm{def:regression\text{-}good\text{-}event}\), at \(d=1\) and \(x=1\) its threshold is
\[
C_m\left(r_n+\sqrt{\frac{1+\log2}{n}}\right).
\]
The definition \(r_n=n^{-\beta/(2\beta+1)}\) and \(\beta>1/2\) give
\[
\frac{\beta}{2\beta+1}>\frac14,
\qquad
r_n+\sqrt{\frac{1+\log2}{n}}
=o\!\left(n^{-1/4}\log n\right).
\]
Thus, for every fixed \(0<C_m<\infty\), the event \(E_n\) is contained in \(\mathcal G_n(1;C_m,\widehat m)^c\) for all sufficiently large \(n\), because the maximum defining \(\mathcal G_n\) includes the selected rotation and the true-direction query just bounded. Therefore
\[
\begin{aligned}
\liminf_{n\to\infty}
P^{\circ\otimes n}\!\left\{\mathcal G_n(1;C_m,\widehat m)^c\right\}
&\ge\lim_{n\to\infty}P^{\circ\otimes n}(E_n)\\
&=e^{-\delta}>2e^{-1}.
\end{aligned}
\]
The same fixed law \(P^\circ\) and the same \(\delta\) work for every finite \(C_m>0\).

It remains to verify the stated protocol-qualified consequence. Suppose that statistics and radii within \(\mathfrak H_{\mathrm{orth}}(\mathcal A_n^{\mathrm{LP}})\) satisfied the displayed uniform probability requirement for some \(c_u,C_H,C_m>0\). At \(d=1\) and \(x=1\), put \(u_1=1+\log2\). Since \(c_u>0\), one has \(u_1\le c_u n\) for all sufficiently large \(n\). If \(\mathcal R_n\) denotes the residual-HSIC error event appearing second in that failure union, set inclusion gives
\[
P^{\circ\otimes n}\!\left(
\mathcal G_n(1;C_m,\widehat m)^c\cup\mathcal R_n
\right)
\ge
P^{\circ\otimes n}\!\left\{\mathcal G_n(1;C_m,\widehat m)^c\right\}.
\]
Taking lower limits makes the left side strictly larger than \(2e^{-1}\), whereas the required inequality bounds it by \(2e^{-1}\) for every sufficiently large \(n\), a contradiction. This contradiction is independent of the form of \(\widehat H^\perp\) and of its radii, but it uses the compulsory event \(\mathcal G_n\) and the fixed class \(\mathcal A_n^{\mathrm{LP}}\). It therefore yields no assertion about regression estimators, independence statistics, or residual-HSIC procedures outside that protocol.
\end{proof}
\par\smallskip\noindent \textit{Justification.} A single admissible nonlinear Hermite-chart law has predictor regions of probability order \(1/n\) that the fixed-bandwidth regression fold misses with limiting probability \(e^{-\delta}\). On such a missed region the ridge-only local-polynomial fit is zero while the quadratic conditional mean is of order \(\log n\), forcing an \(L^4\) error larger than the protocol threshold. Since the exact orthogonalization failure union compulsorily contains the regression-good-event complement, its required \(2/e\) bound at \(x=1\) is impossible. \textit{Gap.} Nuisance-aware residual theory and pointwise consistency for suitable learners do not validate the note's exact uniform four-fold local-polynomial/orthogonalization probability requirement on the unbounded positive-density envelope; the theorem supplies a concrete fixed-law obstruction to that requirement. \textit{Consumer.} Method developers should not use this particular regression-good event to justify uniform residual-HSIC inference under the stated envelope. The result does not rule out a different regression estimator, a different independence statistic, or a different HSIC-error protocol.

\begin{lemma}[lem:subhalf-chart-calibration]\label{lem:subhalf-chart-calibration}
The existential chart constants in Theorem \(\mathrm{thm:attainable\text{-}margin\text{-}rectangle}\) may be selected with \(0<a_-<a_+<1/2\).  For such a selection there are \(t_\star>0\), \(0<\lambda_H\le\lambda_0\), and constants \(0<c_H<C_H<\infty\) such that every \((P_{a,t,\theta},\theta)\in\mathcal M_d^{\mathrm H}(\lambda)\) satisfies, for every \(j\in I\),
\[
\mathbb E_{P_{a,t,\theta}}\psi(X_j,Y_j)=\theta_j\,2t_j(1-2a_j),
\qquad
\operatorname{Var}_{P_{a,t,\theta}}\{\psi(X_j,Y_j)\}\le V_\psi,
\]
and
\[
c_Ht_j^2\le H_{j,-\theta_j}(P_{a,t,\theta})\le C_Ht_j^2.
\]
The class \(\mathcal M_d^{\mathrm H}(\lambda)\) is nonempty for every \(d\ge1\) and \(0<\lambda\le\lambda_H\).  More precisely, there are \(a_0\in(a_-,a_+)\), a continuous choice \(t_\lambda\in(0,t_\star]\), and \(C_I<\infty\) such that \(H_{1,-}(P_{a_0,t_\lambda,+})=\lambda\), both coordinate-swapped orientations belong to \(\mathcal M_1^{\mathrm H}(\lambda)\), and
\[
\operatorname{KL}(P_{a_0,t_\lambda,+}\Vert P_{a_0,t_\lambda,-})\le C_I\lambda.
\]
\end{lemma}
\begin{proof}
Choose, for example, a compact interval \(A=[1/8,1/4]\subset(0,1/2)\).  The chart membership calculation underlying the already-proved Theorem \(\mathrm{thm:attainable\text{-}margin\text{-}rectangle}\) is uniform on every compact subset of \((0,1)\): for \(a\in A\) and sufficiently small positive \(t\),
\[
C=Z,\qquad E=aZ+t(Z^2-1)+s_{a,t}\xi,qquad
s_{a,t}^2=1-a^2-2t^2,
\]
with independent standard normal \(Z,\xi\).  Taking \(t_\star\) small makes \(s_{a,t}^2\) uniformly bounded away from zero.  The cause and noise densities are then everywhere positive Gaussian densities, \(C\perp\xi\), and the product construction gives independent pairs.  Also
\[
\operatorname{Var}(C)=1,qquad
\operatorname{Var}(E)=a^2+2t^2+s_{a,t}^2=1.
\]
The mechanism is quadratic and both log densities are quadratic, so their weighted H\"older norms in \(\mathcal H_{(2)}^\beta(L)\) are uniformly finite after selecting the fixed existential envelope in that theorem.  The same chart calculation verifies Peters et al.'s restricted-ANM Condition 19 for every \(t>0\); the excluded boundary \(t=0\) is exactly the reversible linear-Gaussian member.  Hence all eight properties in Definition \(\mathrm{def:base\text{-}model}\) hold.  The one-sided expansion from Theorem \(\mathrm{thm:attainable\text{-}margin\text{-}rectangle}\) gives, uniformly on \(A\),
\[
H_{j,-\theta_j}=\kappa_H(a_j)t_j^2+O(t_j^3),
\qquad
g_j=\frac{a_j^2}{2}+O(t_j^2).
\]
The Gaussian integral defining \(\kappa_H\) is continuous in \(a\) by dominated convergence, because both kernels are bounded by one and the chart densities have a common Gaussian envelope.  The established positivity of the coefficient therefore yields
\[
0<\kappa_-:=\inf_{a\in A}\kappa_H(a)
\le \sup_{a\in A}\kappa_H(a)=:\kappa_+<\infty.
\]
Decrease \(t_\star\) until the uniform remainder has absolute value at most \(\kappa_-t^2/2\).  Then the asserted HSIC inequalities hold with \(c_H=\kappa_-/2\) and, for example, \(C_H=\kappa_++\kappa_-/2\).  Choose \(\underline g\) strictly below \((1/8)^2/2\) and \(\overline g\) strictly above \((1/4)^2/2\), and decrease \(t_\star\) once more; the gain expansion then gives the navigation window throughout the subchart.

It remains to compute the observable moment.  In the orientation \(\theta_j=+1\), write \(X=Z\) and \(Y=aZ+t(Z^2-1)+s\xi\).  Centering and the standard-normal moments \(\mathbb EZ^2=1\), \(\mathbb EZ^4=3\) give
\[
\mathbb E[(X^2-1)Y]
=t\mathbb E[(Z^2-1)^2]=2t.
\]
Moreover,
\[
\begin{aligned}
\mathbb E[(Y^2-1)X]
&=\mathbb E[ZY^2]\\
&=2at\,\mathbb E[Z^2(Z^2-1)]
=4at,
\end{aligned}
\]
because every remaining term is an odd centered Gaussian moment or contains the independent centered \(\xi\).  Thus \(\mathbb E\psi(X,Y)=2t(1-2a)\).  Coordinate swapping changes \(\psi(x,y)\) to \(-\psi(x,y)\), proving the formula for both orientations.

For the variance, \(|a|,t,s\le1\) implies
\[
|Y|^8\le3^7\{|Z|^8+|Z^2-1|^8+|\xi|^8\},
\qquad \mathbb E|Y|^8\le M_8.
\]
The elementary inequality \((u^2-1)^2\le2(u^4+1)\), followed by Cauchy--Schwarz and Lyapunov's moment inequality, yields
\[
\begin{aligned}
\mathbb E\psi(X,Y)^2
&\le4\{\mathbb E(X^4Y^2)+\mathbb EY^2+\mathbb E(Y^4X^2)+\mathbb EX^2\}\\
&\le4\left\{(\mathbb E|Z|^8)^{1/2}M_8^{1/4}+2
+M_8^{1/2}(\mathbb E|Z|^4)^{1/2}\right\}=V_\psi.
\end{aligned}
\]
Since variance is at most the second moment, the claimed variance bound follows.

Finally take \(a_0=(a_-+a_+)/2=3/16\).  The one-sided Jacobian assertion in Theorem \(\mathrm{thm:attainable\text{-}margin\text{-}rectangle}\), equivalently the expansion with positive coefficient just used, makes \(t^2\mapsto H_{1,-}(P_{a_0,t,+})\) locally invertible at zero.  Hence, after choosing \(\lambda_H>0\) small enough, for each \(0<\lambda\le\lambda_H\) there is \(t_\lambda\le t_\star\) with \(H_{1,-}=\lambda\), and \(t_\lambda^2\asymp\lambda\).  The gain window already proved and \(A_H>1\) put both coordinate-swapped laws in \(\mathcal M_1^{\mathrm H}(\lambda)\).  Proposition \(\mathrm{prop:mirror\text{-}information}\), evaluated at this midpoint, gives
\[
\operatorname{KL}(P_{a_0,t_\lambda,+}\Vert P_{a_0,t_\lambda,-})
\le C_I\lambda.
\]
Taking products of these one-pair members proves nonemptiness for every \(d\).
\end{proof}

\begin{lemma}[lem:chart-product-testing-lower]\label{lem:chart-product-testing-lower}
There is a constant \(c_L>0\), depending only on the fixed sub-half chart, such that for every \(d\ge1\), \(0<\alpha\le1/4\), \(0<\lambda\le\lambda_H\), and every measurable estimator \(\widetilde\theta\),
\[
n\lambda\le c_L\log\frac{d+1}{\alpha}
\quad\Longrightarrow\quad
\sup_{(P,\theta)\in\mathcal M_d^{\mathrm H}(\lambda)}
P^{\otimes n}(\widetilde\theta\ne\theta)>\alpha.
\]
In the same regime, no random set can have both uniform coverage at least \(1-\alpha/2\) and uniform singleton probability at least \(1-\alpha/2\) over \(\mathcal M_d^{\mathrm H}(\lambda)\).
\end{lemma}
\begin{proof}
Let
\[
Q_\lambda^+=P_{a_0,t_\lambda,+},
\qquad
Q_\lambda^-=P_{a_0,t_\lambda,-}
\]
be the coordinate-swapped laws supplied by Lemma \(\mathrm{lem:subhalf\text{-}chart\text{-}calibration}\).  That lemma gives the two typed law--orientation memberships
\[
\bigl(Q_\lambda^+,+1\bigr),
\bigl(Q_\lambda^-,-1\bigr)\in\mathcal M_1^{\mathrm H}(\lambda),
\qquad
\operatorname{KL}(Q_\lambda^+\Vert Q_\lambda^-)
\le C_I\lambda.
\tag{15}
\]
For \(\vartheta=(\vartheta_1,\ldots,\vartheta_d)\in\{-1,+1\}^d\), define
\[
Q_\vartheta=\bigotimes_{j=1}^d Q_\lambda^{\vartheta_j}.
\tag{16}
\]
The coordinatewise chart parameters in (16) are \(a_j=a_0\) and \(t_j=t_\lambda\).  The pair-product assumption \(\mathrm{ass:pair\text{-}product}\), the product construction and coordinatewise membership requirements in Definition \(\mathrm{def:subhalf\text{-}hermite\text{-}class}\), and the one-coordinate membership in (15) show that
\[
(Q_\vartheta,\vartheta)\in\mathcal M_d^{\mathrm H}(\lambda)
\quad\text{for every }\vartheta\in\{-1,+1\}^d.
\tag{17}
\]
Thus (16) is a finite product subexperiment of the class in the claim.

We first establish an elementary binary bound.  For probability measures \(R,S\) with densities \(r,s\) relative to a common dominating measure, let
\[
A(R,S)=\int\sqrt{rs}
\]
be their Hellinger affinity.  Jensen's inequality, applied under \(R\), gives
\[
\operatorname{KL}(R\Vert S)
=-2\mathbb E_R\log\sqrt{s/r}
\ge-2\log A(R,S).
\tag{18}
\]
Also, Cauchy--Schwarz applied to
\(|r-s|=|\sqrt r-\sqrt s|(\sqrt r+\sqrt s)\) gives
\[
\|R-S\|_{\mathrm{TV}}
=\frac12\int|r-s|
\le\sqrt{1-A(R,S)^2}.
\tag{19}
\]
Indeed, the two squared integrals in Cauchy--Schwarz are \(2-2A(R,S)\) and \(2+2A(R,S)\).  The equal-prior binary Bayes error is
\[
e(R,S)=\frac{1-\|R-S\|_{\mathrm{TV}}}{2}.
\]
Combining (18)--(19) therefore proves
\[
e(R,S)\ge
\frac{1-\sqrt{1-\exp\{-\operatorname{KL}(R\Vert S)\}}}{2}.
\tag{20}
\]

Apply (20) to
\(R=(Q_\lambda^+)^{\otimes n}\) and
\(S=(Q_\lambda^-)^{\otimes n}\).  Product factorization of the likelihood ratio and (15) yield
\[
K_n:=\operatorname{KL}(R\Vert S)
=n\operatorname{KL}(Q_\lambda^+\Vert Q_\lambda^-)
\le C_In\lambda.
\tag{21}
\]
Let \(e_n=e(R,S)\).  Under the uniform prior on
\(\vartheta\in\{-1,+1\}^d\), both the prior and the likelihood of the experiment (16) factor over \(j\).  Consequently the posterior factorizes, the Bayes-optimal vector decision takes the coordinatewise posterior maximizers, and its probability of recovering every coordinate is \((1-e_n)^d\).  Hence every estimator \(\widetilde\theta\) obeys
\[
\sup_{\vartheta\in\{-1,+1\}^d}
Q_\vartheta^{\otimes n}(\widetilde\theta\ne\vartheta)
\ge 1-(1-e_n)^d.
\tag{22}
\]

It remains to choose one constant uniformly in \(d\) and \(\alpha\).  Put
\[
\gamma=\frac1{64},
\qquad
R_\alpha=\frac{d+1}{\alpha},
\]
and choose \(c_L>0\) so that \(C_Ic_L<\gamma\).  If
\(n\lambda\le c_L\log R_\alpha\), then (21) gives
\[
e^{-K_n}>R_\alpha^{-\gamma}.
\tag{23}
\]
For \(d\le2\), since \(R_\alpha\le3/\alpha\),
\[
R_\alpha^{-\gamma}
\ge(3/\alpha)^{-\gamma}>4\alpha(1-\alpha),
\qquad 0<\alpha\le1/4.
\tag{24}
\]
To verify the strict inequality in (24), divide its two sides: the resulting ratio is decreasing on \((0,1/4]\), and at \(\alpha=1/4\) it is
\(12^{-1/64}/(3/4)>1\).  Equations (20), (23), and (24) give \(e_n>\alpha\), because
\((1-2\alpha)^2=1-4\alpha(1-\alpha)\).  Thus (22) is strictly greater than \(\alpha\).

For \(d\ge3\), the elementary inequality
\(1-\sqrt{1-x}\ge x/2\) for \(0\le x\le1\), together with (20) and (23), gives
\[
e_n>\frac14R_\alpha^{-\gamma}.
\tag{25}
\]
Set \(x=d/\alpha\).  Since \(x\ge12\) and
\(R_\alpha=(d+1)/\alpha\le(4/3)x\), (25) implies
\[
de_n>
\frac{x\alpha}{4}\{(4/3)x\}^{-\gamma}
\ge2\alpha.
\tag{26}
\]
The last inequality follows from
\(x^{1-\gamma}\ge8(4/3)^\gamma\), which holds at \(x=12\) and whose left side is increasing.  Using \(1-u\le e^{-u}\), if \(de_n\le1\), then
\[
1-(1-e_n)^d\ge1-e^{-de_n}\ge\frac{de_n}{2}>\alpha;
\]
if \(de_n>1\), then the same quantity is greater than
\(1-e^{-1}>1/4\ge\alpha\).  Combining this with (17) and (22) proves the estimator lower bound for every \(d\ge1\).

For the random-set assertion, suppose to the contrary that a measurable random set \(\widetilde{\mathcal C}\) had, under every member of the chart class, both coverage at least \(1-\alpha/2\) and singleton probability at least \(1-\alpha/2\).  Define an estimator to be the unique element of \(\widetilde{\mathcal C}\) on the singleton event and a fixed orientation otherwise.  For every \((P,\theta)\in\mathcal M_d^{\mathrm H}(\lambda)\),
\[
\begin{aligned}
P^{\otimes n}(\widetilde\theta\ne\theta)
&\le P^{\otimes n}\{\theta\notin\widetilde{\mathcal C}\}
+P^{\otimes n}\{|\widetilde{\mathcal C}|\ne1\}\\
&\le\alpha.
\end{aligned}
\]
This contradicts the strict estimator lower bound just proved in the same regime.  Therefore the two uniform random-set guarantees cannot coexist.
\end{proof}

\begin{lemma}[lem:peters-disjoint-pair-specialization]\label{lem:peters-disjoint-pair-specialization}
\texttt{For a directed acyclic graph that is a disjoint union of directed edges,\allowbreak{} the multivariate restricted additive-\allowbreak{}noise condition in Peters et al.'s Definition 27 specializes to the conjunction of their bivariate Condition 19 over the edges;\allowbreak{} under that condition their Theorem 28 identifies the directed acyclic graph from the observational law.}
\end{lemma}

\section*{Technical internal limitation (diagnostic only)}

The unconditional upper frontier and confidence set are chart-specific. Their proof uses the exact identity \(\mathbb E\psi(X_j,Y_j)=\theta_j2t_j(1-2a_j)\), the strict sub-half restriction \(a_j<1/2\), and the calibration \(t_j^2\asymp H_{j,-\theta_j}\); none of these identities is asserted for a general member of \(\mathcal M_d(\lambda)\). Thus the larger restricted-ANM class receives only the transferred lower bound. The local-polynomial theorem has a different scope: its witness violates the exact fixed-protocol probability requirement because the required failure union contains \(\mathcal G_n^c\). It proves neither a residual-HSIC estimation lower bound nor impossibility for arbitrary regression estimators, arbitrary independence statistics, or all HSIC-error procedures.

\section{Honest scope}

The field headline is exactly \(N_{\mathrm H}^*(d,\alpha,\lambda)\asymp\lambda^{-1}\log((d+1)/\alpha)\) on the explicit sub-half Hermite class \(\mathcal M_d^{\mathrm H}(\lambda)\), together with its computable simultaneous confidence set and matching product-testing converse. The attainable rectangle is a population statement for the positive-density restricted-ANM chart. Because \(\mathcal M_d^{\mathrm H}(\lambda)\subset\mathcal M_d(\lambda)\), only the minimax lower bound transfers upward; no full-class rate upper bound, confidence construction, or nuisance-uniform residual-HSIC expansion is claimed. For \(d=1\), the rule reduces to inversion of one chart-moment interval and has frontier \(\lambda^{-1}\log(2/\alpha)\). The fixed local-polynomial result says only that its exact orthogonalization guarantee is impossible when it compulsorily includes the failed event \(\mathcal G_n\); it does not constrain procedures outside that protocol. The oracle residual-HSIC score is consistently described as GENE-inspired rather than as the published GENE score.

\begin{thebibliography}{99}

  \bibitem{pmlr-v267-li25bx} Li, M., Qian, H., Wang, T.-Z., Li, S., Zhang, M., and Zhou, A. (2025). Strong and weak identifiability of optimization-based causal discovery in non-linear additive noise models. Proceedings of Machine Learning Research 267, 35753--35768.

  \bibitem{PetersMJS2014} Peters, J., Mooij, J. M., Janzing, D., and Schoelkopf, B. (2014). Causal discovery with continuous additive noise models. Journal of Machine Learning Research 15, 2009--2053.

  \bibitem{Mooijetal16} Mooij, J. M., Peters, J., Janzing, D., Zscheischler, J., and Schoelkopf, B. (2016). Distinguishing cause from effect using observational data: methods and benchmarks. Journal of Machine Learning Research 17, 1--102.

  \bibitem{GrettonBSS05} Gretton, A., Bousquet, O., Smola, A., and Schoelkopf, B. (2005). Measuring statistical dependence with Hilbert--Schmidt norms. Algorithmic Learning Theory, 63--77.

  \bibitem{Pfister2018dHSIC} Pfister, N., Buehlmann, P., Schoelkopf, B., and Peters, J. (2018). Kernel-based tests for joint independence. Journal of the Royal Statistical Society, Series B 80, 5--31.

  \bibitem{HlavkaHuskovaMeintanis2011} Hl\'avka, Z., Hu\v{s}kov\'a, M., and Meintanis, S. G. (2011). Tests for independence in non-parametric heteroscedastic regression models. Journal of Multivariate Analysis 102, 816--827.

  \bibitem{ShekharKimRamdas2023} Shekhar, S., Kim, I., and Ramdas, A. (2023). A permutation-free kernel independence test. Journal of Machine Learning Research 24, 1--68.

  \bibitem{KimBalakrishnanWasserman2022} Kim, I., Balakrishnan, S., and Wasserman, L. (2022). Minimax optimality of permutation tests. Annals of Statistics 50, 225--251.

  \bibitem{WangKolarDrton2026} Wang, Y. S., Kolar, M., and Drton, M. (2026). Confidence sets for causal orderings. Journal of the American Statistical Association 121(553), 690--703. doi:10.1080/01621459.2025.2542552.

  \bibitem{GaoStoev2020} Gao, Z. and Stoev, S. (2020). Fundamental limits of exact support recovery in high dimensions. Bernoulli 26, 2605--2638.

  \bibitem{Stone1982} Stone, C. J. (1982). Optimal global rates of convergence for nonparametric regression. Annals of Statistics 10, 1040--1053.

  \bibitem{Wainwright2019} Wainwright, M. J. (2019). High-Dimensional Statistics: A Non-Asymptotic Viewpoint. Cambridge University Press.

  \bibitem{Tsybakov2009} Tsybakov, A. B. (2009). Introduction to Nonparametric Estimation. Springer.

  \bibitem{MartingaleHSIC2026} Laumann, F., Liu, Z., and Barahona, M. (2026). A martingale kernel independence test. Public preprint, arXiv:2605.22549.

  \bibitem{AcharyaBhadaneBhattacharyyaKandasamySun2023} Acharya, J., Bhadane, S., Bhattacharyya, A., Kandasamy, S., and Sun, Z. (2023). Sample complexity of distinguishing cause from effect. Proceedings of the 26th International Conference on Artificial Intelligence and Statistics, Proceedings of Machine Learning Research 206, 10487--10504.

  \bibitem{Jin2026LiNGAM} Jin, J. (2026). How many samples are needed to determine causal direction? Sharp minimax bounds for bivariate LiNGAM. Public preprint, arXiv:2608.15840v1, submitted August 16, 2026.

\end{thebibliography}

\end{document}